\chapter{电磁感应与场的表示}
电动势是单位电荷沿回路所受力的线积分，磁通是磁场穿过曲面的面积分。把二者联系起来，需要确定回路的运动方式，并统一回路正向与磁通法向。滑动触点和滚动接触还涉及两种不同的速度：导体材料的速度与几何接触位置的速度。

以下采用SI制。闭合曲线的正向与曲面法向按右手规则对应；$\vct u$表示导体材料速度，$\vct v$表示粒子速度。外加场由给定源维持，若不另说明，忽略感应电流对外加场的反作用。



\section{回路、电动势与磁通的变化}\label{sec:i1-loop}
运动回路的电动势是单位电荷沿导体路径所受力的闭路积分。回路方向和材料速度确定后，磁通变化可分成场的时间变化与边界扫过面积两部分。

\subsection{沿回路积分的是哪一种力}
设载流子相对细导线的漂移沿线切向，导体材料速度为$\vct u$。载流子的完整磁力还含漂移速度项，但该项垂直于漂移方向，在沿线积分中投影为零。因此，单位正电荷沿一瞬间导电路径$C(t)$的电磁驱动定义为
\begin{equation}\label{eq:i1-emf}
\mathcal E=\oint_{C(t)}(\vct E+\vct u\times\vct B)\cdot\dd\vct l.
\end{equation}
若回路含电池，还须另列化学力所作的单位电荷功，作为电源电动势。

固定导线时$\vct u=0$，Maxwell--Faraday方程给出
\[\oint_C\vct E\cdot\dd\vct l
=-\int_S\partial_t\vct B\cdot\dd\vct S.
\]
固定积分面的偏导数只描述场随时间改变。导线运动时，同一瞬间围成的曲面也在改变，不能继续只对$B$求偏导。

\subsection{扫过面积项的符号}
把曲面边缘在$\dd t$内扫出的窄条补到旧曲面上。按规定的边界方向，该窄条的有向面积为$\vct u\times\dd\vct l\,\dd t$，所以
\[\begin{aligned}
\frac{\dd\Phi_B}{\dd t}
&=\int_{S(t)}\partial_t\vct B\cdot\dd\vct S
+\oint_{C(t)}\vct B\cdot(\vct u\times\dd\vct l)\\
&=\int_{S(t)}\partial_t\vct B\cdot\dd\vct S
-\oint_{C(t)}(\vct u\times\vct B)\cdot\dd\vct l.
\end{aligned}\]
第二行用了标量三重积。一个可核对符号的构形是：正$z$磁场中的圆环径向扩大，正向为逆时针。此时$\vct u\times\vct B$沿顺时针，动生积分为负，而磁通增长为正。

将两项相加，便得材料回路的磁通定律
\begin{equation}\label{eq:i1-flux-rule}
\boxed{\mathcal E=-\frac{\dd\Phi_B}{\dd t}}.
\end{equation}
曲面选取不影响结果：两个具有相同边界的曲面合成闭合面，其磁通差由$\nabla\cdot\vct B=0$保证为零。穿过奇异细导线时应先用有限截面的模型，再取极限。

滑动刷接、同极发电机和导线换接需要另辨回路的身份。若画出的几何回路保持不动，其中的导体却持续换成别的材料，几何边界速度就不等于式\eqref{eq:i1-emf}中的材料速度。这时应直接逐段计算洛伦兹力积分；不能以纸面上围住的面积不变为由令电动势为零。滚轮例题中将同时使用直接积分与瞬心方法。


\begin{example}[扩大的圆环与随时间变化的磁场]\label{ex:i1-e1}
半径为$a(t)$的圆形导线环位于$xy$平面，外磁场为$B(t)\vct e_z$。环上每个材料点径向运动，导线电阻为$R$，忽略自感。求感生、动生和总电动势，并在$B>0,\dot B>0,\dot a>0$时确定电流方向。 见图~\ref{fig:i1-e1-q}。
\end{example}

\begin{figure}[H]\centering
\begin{tikzpicture}[>=Stealth,line cap=round,line join=round]
\draw[curve](0,0)circle(1.4);
\draw[axis](-1.9,0)--(2.5,0)node[right]{$x$};\draw[axis](0,-1.8)--(0,2.1)node[above]{$y$};
\node at(-.65,.65){$\odot B(t)$};\draw[->](0,0)--(35:1.4)node[midway,above left]{$a$};
\draw[->](1.4,.2)--(2.15,.2)node[midway,above]{$\dot a$};
\node[below left]at(0,0){$O$};
\end{tikzpicture}
\caption{回路正向取逆时针，磁通正向朝纸外。}\label{fig:i1-e1-q}
\end{figure}
\begin{solution}
取图~\ref{fig:i1-e1-q}的逆时针方向为正，故面积正向为$+z$。固定环形路径求感生部分时，只让磁场随时间变：
\[\mathcal E_{\rm tr}=-\pi a^2\dot B.
\]
动生部分来自每个材料点的径向速度$\vct u=\dot a\vct e_r$。由$\vct e_r\times\vct e_z=-\vct e_\theta$，沿逆时针线元有
\[\mathcal E_{\rm mot}=\int_0^{2\pi}(-B\dot a\vct e_\theta)\cdot(a\vct e_\theta\dd\theta)
=-aB\dot a[\theta]_0^{2\pi}
=-2\pi aB\dot a.
\]
图~\ref{fig:i1-e1-s}左图给出该项对应的窄条面积。它是整个边缘的扫掠，不能只取一个材料点的位移乘磁场。


\begin{figure}[H]\centering
\begin{tikzpicture}[>=Stealth,line cap=round,line join=round]
\begin{scope}[xshift=-2.5cm]\draw[curve](0,0)circle(1.2);\draw[guide](0,0)circle(1.5);\draw[->](1.2,0)--(1.5,0)node[midway,above]{$\dd a$};\draw[->](1.2,0)--(1.2,-.7)node[right]{$\vct u\times\vct B$};\node[below]at(0,-1.7){扫过面积$2\pi a\,\dd a$};\end{scope}
\begin{scope}[xshift=2.2cm]\draw[curve](0,0)circle(1.2);\draw[->](0,1.2)arc(90:-40:1.2);\node at(0,0){$I<0$};\node[below]at(0,-1.7){顺时针感应电流};\end{scope}
\end{tikzpicture}
\caption{动生项的符号由材料速度与磁场的叉积确定。}\label{fig:i1-e1-s}
\end{figure}

合并两部分，
\[\boxed{\mathcal E=-\pi a^2\dot B-2\pi aB\dot a
=-\frac{\dd}{\dd t}(\pi a^2B)},\qquad I=\frac{\mathcal E}{R}.
\]
给定的三个量都为正时，两项同号，$I<0$，即图~\ref{fig:i1-e1-s}右图的顺时针方向。电流产生的磁场朝纸内，阻碍原来朝纸外的磁通增长。若磁场减小而圆环扩大，两项可以抵消；总电动势是否为零取决于$a^2B$是否保持不变，不能只观察其中一种变化。
\end{solution}


\section{局部场方程与微分运算}\label{sec:i1-maxwell}
回路和曲面的积分关系可以缩小到局部，但奇异源与界面须单独保留。下面用同一套方向约定连接积分式、微分式和跨界面的跃变条件。

\subsection{从积分式到微分式}
在固定空间区域，Gauss定理把通量积分化为散度积分，Stokes定理把边界环量化为旋度积分。由于区域可以任意缩小，真空中的四条关系为
\[\nabla\cdot\vct E=\frac\rho{\epsilon_0},\qquad
\nabla\cdot\vct B=0,
\]
\[\nabla\times\vct E=-\partial_t\vct B,\qquad
\nabla\times\vct B=\mu_0\vct j+\mu_0\epsilon_0\partial_t\vct E.
\]
最后一式的位移电流不能在所有问题中略去。对其取散度，左端恒为零，得到
\[0=\mu_0\nabla\cdot\vct j+\mu_0\partial_t\rho,
\]
即电荷连续性方程。若去掉位移电流，却又允许局部电荷密度任意改变，两条方程便不能同时成立。

以边长$\dd x,\dd y,\dd z$的小长方体为例，左右两个面的$B_x$通量之和为$[B_x(x+\dd x)-B_x(x)]\dd y\dd z$。再加其余四面并除以体积，极限给出$\partial_xB_x+\partial_yB_y+\partial_zB_z$。对$xy$小矩形作逆时针环量，同样得到
\[\oint\vct E\cdot\dd\vct l
=(\partial_xE_y-\partial_yE_x)\dd x\dd y+o(\dd x\dd y).
\]
该式对应旋度的$z$分量；换成另外两个坐标面便得到其余分量。

\subsection{两条恒等式及其限制}
令$\epsilon_{ijk}$为Levi-Civita符号，$\partial_i=\partial/\partial x_i$。在二阶偏导可交换的区域，
\[(\nabla\times\nabla f)_i
=\epsilon_{ijk}\partial_j\partial_kf=0,
\]
因为交换$j,k$后偏导不变，而$\epsilon_{ijk}$变号。类似地，
\[\nabla\cdot(\nabla\times\vct A)
=\epsilon_{ijk}\partial_i\partial_jA_k=0.
\]
这两式是微分运算的恒等式，并不依赖材料的性质。它们分别保证梯度场局部无旋、旋度场局部无散。

反命题需要区域条件。局部无旋场在单连通、无奇点的区域可以写成单值标势的梯度；绕过被排除轴线的区域可能有非零环量。磁场无散可在适当区域引入磁矢势，但矢势也可能需要分片定义。处理细长磁通管时，这些全域条件会直接影响回路积分。



\section{电势、磁矢势与规范}\label{sec:i1-potential}
时变磁场中的电场一般需要标量势和磁矢势共同表示。势的选择还取决于规范；同一组场可以对应不同的势，初态和边界条件约束的是物理场。

\subsection{由无散场引入矢势}
在通常的无磁单极区域令
\[\vct B=\nabla\times\vct A.
\]
代入Faraday方程后，
\[\nabla\times(\vct E+\partial_t\vct A)=0.
\]
在可定义单值势的区域，括号中的量可写成$-\nabla\varphi$，所以
\begin{equation}\label{eq:i1-potentials}
\boxed{\vct E=-\partial_t\vct A-\nabla\varphi},\qquad
\vct B=\nabla\times\vct A.
\end{equation}
前一项随磁矢势的建立而改变，后一项包括电荷分布产生的梯度场；两项的划分依赖规范，总电场不依赖。

取任意足够光滑的$\chi(\vct r,t)$，作变换
\[\vct A'=\vct A+\nabla\chi,\qquad
\varphi'=\varphi-\partial_t\chi.
\]
矢势增加的梯度不改变$B$；计算$E'$时，$-\partial_t\nabla\chi$与$+\nabla\partial_t\chi$也恰好相消。因此同一组物理场有许多种势的表示，不能把某个分量的数值当作独立可观测量。

\subsection{静磁矢势的Poisson方程}
静磁条件下$\partial_t\vct E=0$，Ampere方程为$\nabla\times\vct B=\mu_0\vct j$。由
\[\nabla\times(\nabla\times\vct A)
=\nabla(\nabla\cdot\vct A)-\nabla^2\vct A,
\]
再选取库仑规范$\nabla\cdot\vct A=0$，得到
\[\nabla^2\vct A=-\mu_0\vct j.
\]
若电流局限在有限区域，且矢势在远处衰减，可逐分量使用Poisson方程的积分解：
\[\boxed{\vct A(\vct r)=\frac{\mu_0}{4\pi}
\int\frac{\vct j(\vct r')}{|\vct r-\vct r'|}\dd V'}.
\]
无限长导线不满足这里的远区衰减假设，矢势只能在选定参考长度后写成对数形式。磁场本身仍可由它的空间导数唯一确定。


\begin{example}[无限直电流的磁矢势与参考长度]\label{ex:i1-e2}
无限长细直导线沿$z$轴通稳恒电流$I$。构造一个轴向磁矢势，说明为什么需要参考长度$\rho_0$。若电流缓慢变化，改变参考长度时电势应怎样同时改变？ 见图~\ref{fig:i1-e2-q}。
\end{example}

\begin{figure}[H]\centering
\begin{tikzpicture}[>=Stealth,line cap=round,line join=round]
\draw[line width=1pt](0,-1.8)--(0,1.8)node[above]{$I\vct e_z$};\draw[->](0,-1.2)--(0,-.2);
\draw[axis](0,0)--(2.5,0)node[right]{$\rho$};\fill(1.7,0)circle(1.5pt)node[above]{$P$};\draw[<->](0,-.5)--(1.7,-.5)node[midway,below]{$\rho$};
\draw[guide](0,0)ellipse(1.7 and .45);\node[left]at(0,0){$O$};
\end{tikzpicture}
\caption{导线外的矢势可选为轴向，磁场沿圆周。}\label{fig:i1-e2-q}
\end{figure}
\begin{solution}
图~\ref{fig:i1-e2-q}中电流沿$+z$，由对称性可以选$\vct A=A_z(\rho)\vct e_z$。该选择自动满足库仑规范。柱坐标旋度的环向分量为
\[B_\theta=-\frac{\dd A_z}{\dd\rho}.
\]
在导线外，Ampere定律给$B_\theta=\mu_0I/(2\pi\rho)$，因此
\[A_z(\rho)-A_z(\rho_0)
=-\int_{\rho_0}^{\rho}\frac{\mu_0I}{2\pi s}\dd s
=-\frac{\mu_0I}{2\pi}[\ln s]_{\rho_0}^{\rho}
=-\frac{\mu_0I}{2\pi}\ln\frac\rho{\rho_0}.
\]
取$A_z(\rho_0)=0$就得到图~\ref{fig:i1-e2-s}中的对数曲线。图线的斜率而非纵坐标零点决定磁场。


\begin{figure}[H]\centering
\begin{tikzpicture}[>=Stealth,line cap=round,line join=round]
\draw[axis](-.1,0)--(4.2,0)node[right]{$\rho/\rho_0$};\draw[axis](0,-1.4)--(0,1.6)node[above]{$A_z$};
\draw[curve,domain=.3:4,samples=101]plot(\x,{-ln(\x)});\draw[guide](1,-1.3)--(1,1.3);\node[below right]at(1,0){$1$};
\node[align=left,right]at(1.5,1.05){$-\partial_\rho A_z=B_\theta$\\[3pt]改换$\rho_0$只加常数};
\end{tikzpicture}
\caption{对数矢势的斜率确定环向磁场；参考点不改变这个斜率。}\label{fig:i1-e2-s}
\end{figure}

若将参考长度改为$\rho_1$，两种矢势相差$\Delta A_z=\mu_0I\ln(\rho_1/\rho_0)/(2\pi)$。它可写成$\nabla\chi$，其中$\chi=\Delta A_z\,z$；静态下电势不需改变。两表示给出同一$B_\theta$。

若改令$I=I(t)$，这个$\chi$也随时间变。保持观察点$z$不变求偏导，规范变换要求
\[\varphi'-\varphi=-\partial_t\chi
=-\frac{\mu_0\dot I}{2\pi}\ln\frac{\rho_1}{\rho_0}\,z.\]
该电势产生的场为$+\partial_t\Delta A_z\,\vct e_z$，与矢势变化增加的$-\partial_t\Delta A_z\,\vct e_z$相消。只改矢势而把电势强行固定，会改变总电场。在实际变流问题中还须给出回流导线的位置；单独一根无限直导线不能为感生电场提供完整的远区条件。
\end{solution}


\section{带电粒子的能量与正则动量}\label{sec:i1-motion}
机械动量直接描述粒子的运动，正则动量则与所选势及坐标有关。轨道守恒量由时空对称性确定，能量变化还须计入电场、外力和随时间运动的约束所做的功。

\subsection{磁力不做功不等于机械能总守恒}
质点方程为
\[m\dot{\vct v}=q(\vct E+\vct v\times\vct B)+\vct F_{\rm ext}.
\]
点乘$\vct v$，磁力项恒为零，得到
\[\frac{\dd}{\dd t}\frac{mv^2}{2}
=q\vct E\cdot\vct v+\vct F_{\rm ext}\cdot\vct v.
\]
因此纯静磁场中的自由粒子速率不变；若磁场随时间改变并产生电场，粒子仍可从该电场获得能量。电池、约束运动的外力和改变外磁场的电源，均须按各自的功率计入能量交换。

\subsection{拉格朗日量与循环坐标}
非相对论粒子的拉格朗日量为
\[\mathcal L=\frac12m\dot{\vct r}^{\,2}+q\vct A\cdot\dot{\vct r}-q\varphi.
\]
对速度求偏导，正则动量为
\[\vct P=\frac{\partial\mathcal L}{\partial\dot{\vct r}}
=m\vct v+q\vct A.
\]
这个拉格朗日量可逐分量恢复洛伦兹方程。记$\partial_i=\partial/\partial x_i$，重复指标求和；先分别计算
\[\begin{aligned}
\frac{\dd}{\dd t}\frac{\partial\mathcal L}{\partial v_i}
&=m\dot v_i+q\partial_tA_i+qv_j\partial_jA_i,\\
\frac{\partial\mathcal L}{\partial x_i}
&=qv_j\partial_iA_j-q\partial_i\varphi.
\end{aligned}\]
两者相等，移项后为
\[m\dot v_i=q(-\partial_tA_i-\partial_i\varphi)
+qv_j(\partial_iA_j-\partial_jA_i).
\]
第一括号是$E_i$；第二项由$\vct B=\nabla\times\vct A$及置换符号收缩写成$q(\vct v\times\vct B)_i$。因此$\mathcal L$给出的正是$m\dot{\vct v}=q(\vct E+\vct v\times\vct B)$。
若在选定规范下$\mathcal L$不显含$x$，Euler--Lagrange方程给$\dot P_x=0$。这不要求机械动量$mv_x$不变：磁场可以在不同方向间转移机械动量。

规范变换引起的拉格朗日量变化为
\[\delta\mathcal L=q\nabla\chi\cdot\dot{\vct r}+q\partial_t\chi
=q\frac{\dd\chi(\vct r(t),t)}{\dd t}.\]
这个全时间导数不改变运动方程，却使$\vct P$增加$q\nabla\chi$。所以判断某个$P_x$是否守恒时，应检查该规范下的坐标依赖。轨道、速率和导引中心等物理结论应独立于这种表示。由Legendre变换，
\[\begin{aligned}
H&=\vct P\cdot\vct v-\mathcal L\\
&=(m\vct v+q\vct A)\cdot\vct v
-\left(\frac12mv^2+q\vct A\cdot\vct v-q\varphi\right)
=\frac12mv^2+q\varphi.
\end{aligned}\]
势不显含时间时$H$守恒；若势显含时间，则不能先假定它守恒。


\begin{example}[正则动量与匀强磁场中的圆轨道]\label{ex:i1-e3}
质量$m$、电荷$q\ne0$的粒子在匀强磁场$B\vct e_z$中运动，$B>0$，无电场和其他力。用正则动量求平面轨道的导引中心，并比较$\vct A=(-By,0,0)$与$\vct A'=(0,Bx,0)$两种规范。 见图~\ref{fig:i1-e3-q}。
\end{example}

\begin{figure}[H]\centering
\begin{tikzpicture}[>=Stealth,line cap=round,line join=round]
\draw[axis](-2,0)--(2.4,0)node[right]{$x$};\draw[axis](0,-1.7)--(0,2.1)node[above]{$y$};
\draw[curve](.45,.4)circle(1.05);\fill(.45,.4)circle(1.4pt)node[below left]{$C$};\fill(1.5,.4)circle(1.6pt)node[right]{$P$};\draw[->](1.5,.4)--(1.5,-.45)node[right]{$\vct v$};
\node at(-1.3,1.3){$\odot\ B$};
\end{tikzpicture}
\caption{正电荷在正$z$磁场中的回旋方向为顺时针。}\label{fig:i1-e3-q}
\end{figure}
\begin{solution}
取图~\ref{fig:i1-e3-q}的$+z$磁场，并先讨论$q>0$。选Landau规范$\vct A=(-By,0,0)$、$\varphi=0$。拉格朗日量不含$x$，故
\[P_x=mv_x-qBy=\text{常数}.
\]
直接用运动方程也能核对：$m\dot v_x=qBv_y$，因此$\dd(mv_x-qBy)/\dd t=0$。再由$m\dot v_y=-qBv_x$得到另一常数$mv_y+qBx$。


\begin{figure}[H]\centering
\begin{tikzpicture}[>=Stealth,line cap=round,line join=round]
\draw[axis](-1.1,-.5)--(2.9,-.5)node[right]{$x$};
\draw[axis](-.7,-.8)--(-.7,2.2)node[above]{$y$};
\draw[curve](.7,.8)circle(.95);\fill(.7,.8)circle(1.5pt)node[above left]{$(X,Y)$};
\fill(1.65,.8)circle(1.5pt)node[above right]{$P$};
\draw[guide](.7,-.5)--(.7,.8)--(-.7,.8);
\node[below]at(.7,-.5){$X$};\node[left]at(-.7,.8){$Y$};
\draw[->](1.65,.8)--(1.65,.05)node[right]{$v_y$};
\draw[<->](.7,-.18)--(1.65,-.18);
\node[below]at(1.17,-.85){$x-X=-mv_y/(qB)$};
\end{tikzpicture}
\caption{两个机械量的组合给出不随规范改变的导引中心。}\label{fig:i1-e3-s}
\end{figure}

将这两个常数改写成长度，图~\ref{fig:i1-e3-s}中的导引中心为
\[\boxed{X=x+\frac{mv_y}{qB},\qquad
Y=y-\frac{mv_x}{qB}}.
\]
逐项求导分别得到$\dot X=v_x-v_x=0$、$\dot Y=v_y-v_y=0$。由此
\[(x-X)^2+(y-Y)^2=\frac{m^2(v_x^2+v_y^2)}{q^2B^2}.
\]
磁力不做功，右端为常数，故轨道为圆，半径$r_L=mv_\perp/(|q|B)$。图中$P$位于圆右端，$q>0$时速度向下，和洛伦兹力指向圆心一致。

也可选$\vct A'=(0,Bx,0)$。两规范之差为$\nabla(Bxy)$，在新规范下$y$为循环坐标，直接守恒的是$P_y'=mv_y+qBx$。不能要求两规范中的同名正则分量数值相同；上述$X,Y$由位置与机械速度组成，其物理含义不变。若$q<0$，公式中的带符号$q$照用，回旋方向反转，半径仍取正值。
\end{solution}


\section{感生电场的空间分布}\label{sec:i1-symmetry}
闭路环量只给出电场的部分信息。对称性和边界条件限定场的形式，法拉第定律再确定相应系数；开放路径的电压须连同测量回路说明。

\subsection{环量给出的信息}
Faraday方程首先给出旋度或闭路环量，通常不能单独确定电场。若$\vct E_1$满足方程，$\vct E_1-\nabla f$仍有相同旋度。要确定梯度部分，还需电荷分布、导体边界或远处的场。

轴对称问题常能化为一维积分，是因为对称性把沿圆周的电场化为同一个未知量。只给出“全空间的$\dot B$为常量”，却不说明磁场源和附加电荷，不能自动推出某一转动对称的电场。例如叠加任意均匀电场，原有旋度保持不变。

以下平面模型采用明确的共同条件：系统沿$z$方向平移不变，给定$B_z(x,y,t)$的变化，采用准静态近似；没有额外电荷场，$E_z=0$，平面内$\nabla\cdot\vct E=0$。有限磁区还要求电场在远处趋零。有限区域磁场的边界电流由外源维持。远区条件与无额外电荷场的约定排除了可任意叠加的调和梯度场，因而确定了这里的感生电场。

\subsection{面积元叠加与平面旋流核}
令$\alpha(x',y')=-\partial_tB_z(x',y')$。面积元$\dd A'$在其外部贡献的环量为$\alpha\dd A'$。将它视为二维的点旋源，得到
\begin{equation}\label{eq:i1-area-kernel}
\vct E(\vct r)=\frac1{2\pi}\int
\alpha(\vct r')\frac{\vct e_z\times(\vct r-\vct r')}{|\vct r-\vct r'|^2}\dd A'.
\end{equation}
距面积元$\rho$处，分子大小为$\rho$、分母为$\rho^2$，所以核整体按$1/\rho$衰减；当$\alpha>0$时，叉积给出逆时针方向。系数由Faraday方程确定，不含$\mu_0$。这个核与无限直电流的平面磁场核有相同的空间形式，所表示的源则是磁通变化率。

求柱坐标旋度时，应同时计入弧长随半径的改变。半径$\rho$与$\rho+\dd\rho$两条圆弧的环量差为$\partial_\rho(\rho E_\theta)\dd\rho\dd\theta$；两条径向边的差为$-\partial_\theta E_\rho\dd\rho\dd\theta$。除以面积$\rho\dd\rho\dd\theta$，即
\[(\nabla\times\vct E)_z
=\frac1\rho\partial_\rho(\rho E_\theta)-\frac1\rho\partial_\theta E_\rho.
\]
缺少第一个括号中的$\rho$，会把环向电场的旋度算错。


\begin{example}[有限圆磁区的内外感生场]\label{ex:i1-e4}
圆柱形磁区的横截面为半径$R$的圆，区内$\partial_tB_z=-\alpha$，$\alpha>0$，区外磁场不随时间变。按本节约定求平面内各点的感生电场，并分别用环量与旋度核对。 见图~\ref{fig:i1-e4-q}。
\end{example}

\begin{figure}[H]\centering
\begin{tikzpicture}[>=Stealth,line cap=round,line join=round]
\fill[gray!8](0,0)circle(1.45);\draw[curve](0,0)circle(1.45);
\draw[axis](-1.9,0)--(2.1,0)node[right]{$x$};\draw[axis](0,-1.8)--(0,1.9)node[above]{$y$};
\node[below left]at(0,0){$O$};\draw[->](0,0)--(40:1.45)node[midway,above left]{$R$};
\node at(-.65,.65){$\dot B_z$};\fill(1,0)circle(1.5pt)node[below right]{$P$};\draw[guide](0,0)circle(.85);\draw[guide](0,0)circle(1.8);\node[right]at(1.1,-1.15){$\dot B_z=-\alpha$};
\end{tikzpicture}
\caption{虚线圆为两种观察路径，分别位于磁区内部和外部。}\label{fig:i1-e4-q}
\end{figure}
\begin{solution}
在图~\ref{fig:i1-e4-q}中，旋转对称与无附加电荷场共同给出$\vct E=E_\theta(r)\vct e_\theta$。对逆时针同心圆积分时，圆周上$E_\theta$处处相同，故左端为$2\pi rE_\theta$。

若$r<R$，被包围的变化磁通只占磁区的一部分，
\[2\pi rE_\theta=\alpha\pi r^2,\qquad
E_\theta=\frac{\alpha r}{2}.
\]
若$r>R$，包围的变化磁通已经达到全部值，不能再把积分面积写成$\pi r^2$，而应写
\[2\pi rE_\theta=\alpha\pi R^2,\qquad
E_\theta=\frac{\alpha R^2}{2r}.
\]

\begin{figure}[H]\centering
\begin{tikzpicture}[>=Stealth,line cap=round,line join=round]
\draw[axis](-.1,0)--(4.5,0)node[right]{$r/R$};\draw[axis](0,-.1)--(0,2.5)node[above]{$E_\theta/(\alpha R/2)$};
\draw[curve](0,0)--(1,1);\draw[curve,domain=1:4,samples=71]plot(\x,{1/\x});
\draw[guide](1,0)--(1,1);\node[below]at(1,0){$1$};\node[above right]at(1,1){边界连续};
\end{tikzpicture}
\caption{区内电场与半径成正比，区外按半径倒数衰减。}\label{fig:i1-e4-s}
\end{figure}

图~\ref{fig:i1-e4-s}将两个区域接在$r=R$处，两式都给$\alpha R/2$，切向电场连续。电场在原点为零，磁区外虽有$\partial_tB=0$，却仍有非零环量。

用柱坐标旋度公式复算，唯一非零分量为
\[(\nabla\times\vct E)_z=\frac1r\frac{\dd(rE_\theta)}{\dd r}
=\begin{cases}
\displaystyle\frac1r\frac{\dd(\alpha r^2/2)}{\dd r}=\alpha,&r<R,\\[4pt]
\displaystyle\frac1r\frac{\dd(\alpha R^2/2)}{\dd r}=0,&r>R.
\end{cases}\]
外区绕磁区的闭合曲线不能缩到一点，其环量由被排除区域的磁通变化确定。沿完整外圆作一周的单位电荷功为$\alpha\pi R^2$，所以不能在整个外区用单值静电势代替这个感生场。
\end{solution}

\begin{example}[长条、正方形与正三角形中的中心场]\label{ex:i1-e5}
三个有限磁区分别为长宽为$D,H$的矩形、正方形和正三角形，中心均为$O$，区内$\partial_tB_z=-\alpha$。求矩形中心场的线性项，说明长条极限；再说明正方形、正三角形近中心场为何具有相同的线性项，以及这个结论的适用范围。 见图~\ref{fig:i1-e5-q}。
\end{example}

\begin{figure}[H]\centering
\begin{tikzpicture}[>=Stealth,line cap=round,line join=round]
\begin{scope}[xshift=-2.6cm]\draw[curve](-2,-.7)rectangle(2,.7);\draw[axis](-2.2,0)--(2.3,0)node[right]{$x$};\draw[axis](0,-1)--(0,1.35)node[above]{$y$};\node at(-1,.35){$-\dot B_z=\alpha$};\node[below]at(0,-.8){长$D$，宽$H$};\end{scope}
\begin{scope}[xshift=2.7cm]\draw[curve](-1.1,-1.1)rectangle(1.1,1.1);\draw[guide](-1.1,-1.1)--(1.1,1.1);\draw[guide](-1.1,1.1)--(1.1,-1.1);\node[below]at(0,-1.3){正方形};\end{scope}
\end{tikzpicture}
\caption{矩形的长宽比改变两个线性系数，方形和三角形另有旋转对称。}\label{fig:i1-e5-q}
\end{figure}
\begin{solution}
图~\ref{fig:i1-e5-q}中矩形以原点为中心。两条反射对称轴使中心电场为零，线性项只能写成$E_x=-a y$、$E_y=b x$；其旋度给$a+b=\alpha$。系数各占多少仍由矩形的长宽比决定。

由面积核直接求$a$。记$D$为沿$x$方向的全长，$H$为全宽，并令
\[F=\frac{y-y'}{(x-x')^2+(y-y')^2},\qquad
E_x=-\frac\alpha{2\pi}\int_{-D/2}^{D/2}\int_{-H/2}^{H/2}F\,\dd y'\dd x'.\]
核在源点处按对称极限解释。平移变量后求观察点导数等价于对磁区边界求导；利用$\partial_yF=-\partial_{y'}F$，得到
\[\left.\partial_yE_x\right|_O
=\frac\alpha{2\pi}\int_{-D/2}^{D/2}
\left[\frac{-y'}{x'^2+y'^2}\right]_{-H/2}^{H/2}\dd x'.\]
上下边界的两项同号，因此
\[\begin{aligned}
-a&=\left.\partial_yE_x\right|_O
=-\frac{\alpha H}{2\pi}\int_{-D/2}^{D/2}
\frac{\dd x'}{x'^2+H^2/4}\\
&=-\frac{\alpha H}{2\pi}\left[\frac2H\arctan\frac{2x'}H\right]_{-D/2}^{D/2}\\
&=-\frac{2\alpha}{\pi}\arctan\frac DH.
\end{aligned}\]
故中心附近
\[E_x=-\frac{2\alpha}{\pi}\arctan\frac DH\,y+O(r^3),\qquad
E_y=\frac{2\alpha}{\pi}\arctan\frac HD\,x+O(r^3).
\]

\begin{figure}[H]\centering
\begin{tikzpicture}[>=Stealth,line cap=round,line join=round]
\begin{scope}[xshift=-2.7cm]\draw[curve](-2,-.75)rectangle(2,.75);\foreach\y in{-.5,0,.5}{\draw[->](.1,\y)--({.1-1.5*\y},\y);}\node[below]at(0,-1.1){窄条：$E_x\simeq-\alpha y$};\end{scope}
\begin{scope}[xshift=2.8cm]\draw[curve](0,1.5)--(-1.3,-.75)--(1.3,-.75)--cycle;\draw[->](.65,0)arc(0:90:.65);\draw[->](-.65,0)arc(180:270:.65);\fill(0,0)circle(1.4pt)node[below right]{$O$};\node[below]at(0,-1.1){三重对称：中心线性项为旋流};\end{scope}
\end{tikzpicture}
\caption{窄条与高阶旋转对称磁区的局部场不同，不能混用全域结论。}\label{fig:i1-e5-s}
\end{figure}

当$D/H\to\infty$时，便得到图~\ref{fig:i1-e5-s}左图的$E_x=-\alpha y,E_y=0$。这是长条中部极限，不是任意有限矩形的全域解。正方形取$D=H$，两系数同为$\alpha/2$。

对正三角形或正方形，绕中心旋转$2\pi/n$后，线性映射矩阵必须保持不变。$n\ge3$时它只能由单位矩阵与平面转动矩阵组成；散度为零消去单位矩阵部分，旋度为$\alpha$给出
\[\boxed{\vct E=\frac\alpha2\vct e_z\times\vct r+\text{高阶项}}.
\]
正三角形允许二次修正，大小为$O(\alpha r^2/H)$；正方形还具有中心反演对称性，首个修正为$O(\alpha r^3/H^2)$。图~\ref{fig:i1-e5-s}右图的圆弧只表示近中心的线性项，靠近边界时不能据此认定整条电场线为圆。
\end{solution}

\begin{example}[用闭合辅助路径求开放线积分]\label{ex:i1-e6}
正方形磁区内$\partial_tB_z=-\alpha$，中心为$O$。取$A=(0,a)$、$B=(-a,a)$，两点和三角形$OAB$都在磁区内。求沿直线$A\to B$的感生电场线积分，再与沿$A\to O\to B$的结果比较。 见图~\ref{fig:i1-e6-q}。
\end{example}

\begin{figure}[H]\centering
\begin{tikzpicture}[>=Stealth,line cap=round,line join=round]
\draw[curve](-1.65,-1.65)rectangle(1.65,1.65);\coordinate(O)at(0,0);\coordinate(A)at(0,1.2);\coordinate(B)at(-1.2,1.2);
\draw[guide](O)--(A)--(B)--cycle;\draw[->,thick](A)--(B);\fill(O)circle(1.4pt)node[below right]{$O$};\fill(A)circle(1.4pt)node[right]{$A$};\fill(B)circle(1.4pt)node[left]{$B$};
\draw[<->](.4,0)--(.4,1.2)node[midway,right]{$a$};\node at(.85,-.8){$-\dot B_z=\alpha$};
\end{tikzpicture}
\caption{借助两条径向对称线闭合路径，计算直线段上的积分。}\label{fig:i1-e6-q}
\end{figure}
\begin{solution}
图~\ref{fig:i1-e6-q}的$OA$位于正方形对称轴上，$OB$位于对角线上。对没有附加电荷场的感生解，由对称性可知这两条线上电场均无径向分量，因此
\[\int_O^A\vct E\cdot\dd\vct l=0,\qquad
\int_B^O\vct E\cdot\dd\vct l=0.
\]
这两个结论对完整方形解成立，不需要把近中心的线性式延伸到边界。图~\ref{fig:i1-e6-s}左图的三角形按$O\to A\to B\to O$定向，法向为$+z$，面积为$a^2/2$。


\begin{figure}[H]\centering
\begin{tikzpicture}[>=Stealth,line cap=round,line join=round]
\begin{scope}[xshift=-2.2cm]\fill[gray!12](0,0)--(0,1.5)--(-1.5,1.5)--cycle;\draw[curve](0,0)--(0,1.5)--(-1.5,1.5)--cycle;\draw[->](0,1.5)--(-1.1,1.5);\node[below right]at(0,0){$O$};\node[right]at(0,1.5){$A$};\node[left]at(-1.5,1.5){$B$};\node[below]at(-.6,-.35){$S_{OAB}=a^2/2$};\end{scope}
\begin{scope}[xshift=2.4cm]\draw[axis](-1.5,0)--(1.5,0);\draw[line width=.4pt](0,-.4)--(0,0);\draw[->](0,.8)--(-.6,.8)node[left]{$\vct E$};\draw[guide](0,0)--(-1.2,1.2);\draw[->](-.6,.6)--(-1,.2);\node[below]at(0,-.55){轴线与对角线上的径向分量为零};\end{scope}
\end{tikzpicture}
\caption{灰色三角形决定两条开放路径的积分差。}\label{fig:i1-e6-s}
\end{figure}

对这个闭合三角形用Faraday定律，
\[\int_O^A+\int_A^B+\int_B^O
=-\int_{OAB}\partial_tB_z\dd S
=\frac{\alpha a^2}{2},
\]
从而沿上方直线段$A\to B$有
\[\boxed{\mathcal E_{AB}^{\rm straight}=\frac{\alpha a^2}{2}}.
\]
若改沿$A\to O\to B$，积分却为零。两条开放路径的差正好等于它们围成的变化磁通，不能把二者都写成同一个不依赖路径的$\varphi(A)-\varphi(B)$。

实际接入电压表时，表线与被测支路共同组成闭合路径。更改表线的布置会改变感应项；导线中还会建立电荷场以满足导电边界。这里算的是给定外场的开放线积分，并未把它等同于忽略引线几何的静电端电压。
\end{solution}

\begin{example}[两圆磁区重叠产生的匀强电场]\label{ex:i1-e7}
两个半径均为$R$的圆磁区，圆心为$C_1=(-c,0)$和$C_2=(c,0)$，$0<c<R$。左圆的$\partial_tB_z=\alpha$，右圆的$\partial_tB_z=-\alpha$。求重叠区域的感生电场，说明该区能否有局部电势。 见图~\ref{fig:i1-e7-q}。
\end{example}

\begin{figure}[H]\centering
\begin{tikzpicture}[>=Stealth,line cap=round,line join=round]
\draw[curve](-.7,0)circle(1.4);\draw[curve](.7,0)circle(1.4);\draw[axis](-2.5,0)--(2.6,0)node[right]{$x$};\draw[axis](0,-1.7)--(0,1.8)node[above]{$y$};
\fill(-.7,0)circle(1.4pt)node[below]{$C_1$};\fill(.7,0)circle(1.4pt)node[below]{$C_2$};\fill(0,.6)circle(1.4pt)node[right]{$P$};
\node at(-1.25,.75){$\dot B_z=\alpha$};\node at(1.25,.75){$\dot B_z=-\alpha$};
\end{tikzpicture}
\caption{两个场按各自圆心求得，再作矢量相加。}\label{fig:i1-e7-q}
\end{figure}
\begin{solution}
在图~\ref{fig:i1-e7-q}的重叠域，点$P$同时位于两个圆的内部。先分别使用圆磁区内部的场，再相加：
\[\vct E_1=-\frac\alpha2\vct e_z\times(\vct r-\vct C_1),\qquad
\vct E_2=+\frac\alpha2\vct e_z\times(\vct r-\vct C_2).
\]
位置矢量$\vct r$的两项恰好抵消，故
\[\vct E=\frac\alpha2\vct e_z\times(\vct C_1-\vct C_2)
=\boxed{-\alpha c\vct e_y}.
\]

\begin{figure}[H]\centering
\begin{tikzpicture}[>=Stealth,line cap=round,line join=round]
\draw[curve](-.75,0)circle(1.35);\draw[curve](.75,0)circle(1.35);\foreach\y in{-.55,0,.55}{\draw[->](0,\y+.2)--(0,\y-.25);}
\node[right]at(.2,0){$\vct E=-\alpha c\vct e_y$};\draw[<->](-.75,-1.7)--(.75,-1.7)node[midway,below]{$2c$};\node[above]at(0,1.7){重叠域内磁场变化率相消};
\end{tikzpicture}
\caption{重叠域的变化磁通密度为零，仍可有由周围磁区决定的匀强电场。}\label{fig:i1-e7-s}
\end{figure}

图~\ref{fig:i1-e7-s}中的箭头长度和方向相同；这个匀强结论在整个重叠透镜内成立，不只是中心附近近似。

重叠域内$\partial_tB_z=0$，所以局部旋度确为零；匀强场并不与此矛盾。该域可以引入局部标势$\varphi_{\rm loc}=\alpha cy+\text{常数}$。但离开重叠域后，总场不再等于同一个常矢量，不能把这个局部势延伸到整个平面。

若点位于一个圆内、另一个圆外，外侧那一项要改为$\mp\alpha R^2\vct e_z\times(\vct r-\vct C_i)/(2|\vct r-\vct C_i|^2)$；两圆外则两项都用外区式。原来的圆心位移抵消只在两项都为线性函数时成立。
\end{solution}

\begin{example}[旋转的非轴对称磁场]\label{ex:i1-e8}
在$r<R$的圆柱形区域中，$B_z(r,\theta,t)=B_0r\cos(\theta-\omega t)$，区外$\partial_tB_z=0$；$B_0$的单位为磁感应强度除以长度。按本节的平面准静态条件，求全空间感生电场，并说明轴对称积分法为何不能直接使用。 见图~\ref{fig:i1-e8-q}。
\end{example}

\begin{figure}[H]\centering
\begin{tikzpicture}[>=Stealth,line cap=round,line join=round]
\draw[curve](0,0)circle(1.5);\draw[axis](-2,0)--(2.1,0)node[right]{$x$};\draw[axis](0,-1.8)--(0,1.9)node[above]{$y$};
\foreach\x in{.4,.9}{\foreach\y in{-.6,0,.6}{\node at(\x,\y){$\odot$};\node at(-\x,\y){$\otimes$};}}
\draw[->](1.75,0)arc(0:65:1.75)node[above]{$\omega$};\node[below]at(0,-2.1){$t=0:\ B_z=B_0x,\quad r<R$};
\end{tikzpicture}
\caption{圆形边界不代表场幅值轴对称，正负磁区绕中心转动。}\label{fig:i1-e8-q}
\end{figure}
\begin{solution}
图~\ref{fig:i1-e8-q}的磁区形状为圆，但磁场幅值不具轴对称性。令$\xi=\theta-\omega t$，直接求导得
\[\partial_tB_z=B_0\omega r\sin\xi.
\]
场的变化率随角度改变，圆周上的电场也不再处处相等，因而不能将环量写成$2\pi rE_\theta=-\dot\Phi$来求单一的$E_\theta$。圆周总磁通为零只约束电场环量，各点电场仍可非零。

引入流函数$\Psi$，令$\vct E=\vct e_z\times\nabla\Psi$。这样散度自动为零，且
\[E_r=-\frac1r\partial_\theta\Psi,\quad E_\theta=\partial_r\Psi,
\qquad \nabla^2\Psi=-B_0\omega r\sin\xi\quad(r<R).
\]
记$K=B_0\omega$。对一般径向函数$f(r)$，$\partial_\theta^2\sin\xi=-\sin\xi$，所以径向方程为$f''+f'/r-f/r^2=-Kr$。齐次试解$r^s$给$s^2-1=0$，通解为$Ar+C/r$。特解试取$Dr^3$，代入左端为$(6D+3D-D)r=8Dr$，故$D=-K/8$。内解排除原点奇点，外部齐次解排除增长项，于是
\[\Psi_{\rm in}=\left(Ar-\frac{Kr^3}{8}\right)\sin\xi,
\qquad \Psi_{\rm out}=\frac Cr\sin\xi.
\]
界面无额外线电荷时，电场的两个分量连续；对这里的$\sin\xi$角谐波，等价地令$\Psi$及$\partial_r\Psi$在$R$处匹配。两式为
\[AR-\frac{KR^3}{8}=\frac CR,\qquad
A-\frac{3KR^2}{8}=-\frac C{R^2}.
\]
第一式除以$R$得$A-KR^2/8=C/R^2$，与第二式相加，得$2A-KR^2/2=0$。于是$A=KR^2/4$，再回代得$C=KR^4/8$。求分量时，内区的两次偏导分别为
\[\partial_\theta\Psi_{\rm in}=\left(Ar-\frac{Kr^3}{8}\right)\cos\xi,
\qquad \partial_r\Psi_{\rm in}=\left(A-\frac{3Kr^2}{8}\right)\sin\xi;\]
外区则为$\partial_\theta\Psi_{\rm out}=C\cos\xi/r$、$\partial_r\Psi_{\rm out}=-C\sin\xi/r^2$。代入两个分量的定义，得到
\[\boxed{\begin{aligned}
E_r&=-\frac K8(2R^2-r^2)\cos\xi,&
E_\theta&=\frac K8(2R^2-3r^2)\sin\xi\quad(r<R),\\
E_r&=-\frac{KR^4}{8r^2}\cos\xi,&
E_\theta&=-\frac{KR^4}{8r^2}\sin\xi\quad(r>R).
\end{aligned}}
\]

\begin{figure}[H]\centering
\begin{tikzpicture}[>=Stealth,line cap=round,line join=round]
\draw[guide](0,0)circle(1.3);
\draw[->](-1.650,-1.650)--(-1.650,-1.722);
\draw[->](-1.650,-1.100)--(-1.688,-1.192);
\draw[->](-1.650,-0.550)--(-1.754,-0.628);
\draw[->](-1.650,0.000)--(-1.794,0.000);
\draw[->](-1.650,0.550)--(-1.754,0.628);
\draw[->](-1.650,1.100)--(-1.688,1.192);
\draw[->](-1.650,1.650)--(-1.650,1.722);
\draw[->](-1.100,-1.650)--(-1.062,-1.742);
\draw[->](-1.100,-1.100)--(-1.100,-1.262);
\draw[->](-1.100,-0.550)--(-1.274,-0.716);
\draw[->](-1.100,0.000)--(-1.398,0.000);
\draw[->](-1.100,0.550)--(-1.274,0.716);
\draw[->](-1.100,1.100)--(-1.100,1.262);
\draw[->](-1.100,1.650)--(-1.062,1.742);
\draw[->](-0.550,-1.650)--(-0.446,-1.728);
\draw[->](-0.550,-1.100)--(-0.474,-1.266);
\draw[->](-0.550,-0.550)--(-0.848,-0.633);
\draw[->](-0.550,0.000)--(-0.973,0.000);
\draw[->](-0.550,0.550)--(-0.848,0.633);
\draw[->](-0.550,1.100)--(-0.474,1.266);
\draw[->](-0.550,1.650)--(-0.446,1.728);
\draw[->](0.000,-1.650)--(0.144,-1.650);
\draw[->](0.000,-1.100)--(0.034,-1.100);
\draw[->](0.000,-0.550)--(-0.340,-0.550);
\draw[->](0.000,0.000)--(-0.465,0.000);
\draw[->](0.000,0.550)--(-0.340,0.550);
\draw[->](0.000,1.100)--(0.034,1.100);
\draw[->](0.000,1.650)--(0.144,1.650);
\draw[->](0.550,-1.650)--(0.654,-1.572);
\draw[->](0.550,-1.100)--(0.626,-0.934);
\draw[->](0.550,-0.550)--(0.252,-0.467);
\draw[->](0.550,0.000)--(0.127,0.000);
\draw[->](0.550,0.550)--(0.252,0.467);
\draw[->](0.550,1.100)--(0.626,0.934);
\draw[->](0.550,1.650)--(0.654,1.572);
\draw[->](1.100,-1.650)--(1.138,-1.558);
\draw[->](1.100,-1.100)--(1.100,-0.938);
\draw[->](1.100,-0.550)--(0.926,-0.384);
\draw[->](1.100,0.000)--(0.802,0.000);
\draw[->](1.100,0.550)--(0.926,0.384);
\draw[->](1.100,1.100)--(1.100,0.938);
\draw[->](1.100,1.650)--(1.138,1.558);
\draw[->](1.650,-1.650)--(1.650,-1.578);
\draw[->](1.650,-1.100)--(1.612,-1.008);
\draw[->](1.650,-0.550)--(1.546,-0.472);
\draw[->](1.650,0.000)--(1.506,0.000);
\draw[->](1.650,0.550)--(1.546,0.472);
\draw[->](1.650,1.100)--(1.612,1.008);
\draw[->](1.650,1.650)--(1.650,1.578);\fill(0,0)circle(1.4pt)node[below right]{$O$};\node[above]at(0,2){$t=0$的电场方向与相对大小};
\end{tikzpicture}
\caption{箭头由解析场逐点计算，不把圆周平均值当作局部场。}\label{fig:i1-e8-s}
\end{figure}

图~\ref{fig:i1-e8-s}按上述公式在$t=0$逐点画出矢量。中心场为$-KR^2\vct e_x/4$，所以中心磁场为零并不推出中心电场为零。区外虽局部无旋，电场仍按$r^{-2}$衰减。两边在$r=R$逐项相同，且各圆周的$\sin\xi$平均为零，与总环量为零相符。
\end{solution}


\section{刚性导体的运动与无阻回路}\label{sec:i1-motional}
转动导体各点的材料速度不同，动生积分需要沿路径使用局部速度。闭合回路的电流还会产生自磁通；电阻为零时，回路方程约束的是总磁链。

\subsection{平移、转动与瞬心}
刚体上材料点相对选定点$C$的位置为$\vct r$，其速度为$\vct u=\vct V_C+\vct\omega\times\vct r$。均匀磁场中的开放路径积分可以分成平移项与转动项，但两项必须采用相同的端点、方向和参考点。

当$\vct\omega$平行于$\vct B$时，利用三重积展开，
\[(\vct\omega\times\vct r)\times\vct B
=\vct r(\vct\omega\cdot\vct B)-\vct\omega(\vct r\cdot\vct B).
\]
在垂直于转轴的平面内，第二项为零。取$\vct\omega=\omega\vct e_z$、$\vct B=B\vct e_z$，其中$\omega,B$均带符号。由于$\dd(r^2)=2\vct r\cdot\dd\vct r$，同一时刻从端点$A$到$B_1$的积分为
\[\begin{aligned}
\int_A^{B_1}(\vct u\times\vct B)\cdot\dd\vct l
&=(\vct V_C\times\vct B)\cdot(\vct r_{B_1}-\vct r_A)
+\frac{\omega B}{2}\int_{r_A^2}^{r_{B_1}^2}\dd s\\
&=(\vct V_C\times\vct B)\cdot(\vct r_{B_1}-\vct r_A)
+\frac{\omega B}{2}(r_{B_1}^2-r_A^2).
\end{aligned}\]
转动项只取决于端点到转轴的距离，弯曲导线与直线参考路径给出相同结果。若转轴不平行于磁场，不能不加条件地沿用这个端点式。

纯滚动的瞬心指刚体在给定时刻速度为零的位置，通常随运动改变。用它计算某一时刻的速度或电动势是合法的；若要求加速度，仍须回到刚体的平动、转动关系，不能把瞬心当作永久固定点求两次导数。


\begin{example}[双半径滚轮的三种算法]\label{ex:i1-e9}
内外半径分别为$a,b$的同心刚性导体轮在水平面上向右纯滚动，$0<a<b$，角速度大小为$\omega$。取外圆最低点$A$和内圆最高点$B$，磁场均匀、垂直纸面向内，大小为$B$。求沿轮内任意导电路径由$A$到$B$的动生电动势，并用速度分解、瞬心和扫过面积相互验证。 见图~\ref{fig:i1-e9-q}。
\end{example}

\begin{figure}[H]\centering
\begin{tikzpicture}[>=Stealth,line cap=round,line join=round]
\draw[curve](0,0)circle(1.6);\draw[curve](0,0)circle(.8);\draw[curve](-2.2,-1.6)--(2.5,-1.6);
\fill(0,-1.6)circle(1.6pt)node[below]{$A$};\fill(0,.8)circle(1.6pt)node[left]{$B$};\node[below left]at(0,0){$C$};
\draw[->](0,0)--(.9,0)node[below]{$V_C=\omega b$};\draw[->](0,.8)--(1.35,.8)node[above]{$u_B=\omega(a+b)$};
\draw[->](1.9,.6)arc(20:-55:1.9)node[right]{$\omega$};\node at(-1,-.45){$\otimes\,\vct B$};
\draw[guide](0,0)--(145:1.6)node[midway,above]{$b$};\draw[guide](0,0)--(55:.8)node[midway,right]{$a$};
\end{tikzpicture}
\caption{A是外圆瞬时接地点，B在内圆上；电刷只在指定端点接触。}\label{fig:i1-e9-q}
\end{figure}
\begin{solution}
图~\ref{fig:i1-e9-q}中，外圆下端$A$与地面纯滚动接触。取$x$向右、$y$向上，磁场为$-B\vct e_z$，角速度为$-\omega\vct e_z$。轮心速度$\vct V_C=\omega b\vct e_x$，任意点的速度为
\[\vct u=(\omega b+\omega y)\vct e_x-\omega x\vct e_y.
\]
因此
\[\vct u\times\vct B
=B\omega\left[x\vct e_x+(b+y)\vct e_y\right]
=\nabla\frac{B\omega}{2}\left[x^2+(b+y)^2\right].
\]
任意轮内路径$A\to B$的积分相同。两端分别为$(0,-b)$和$(0,a)$，所以
\[\boxed{\mathcal E_{A\to B}=\frac{B\omega}{2}(a+b)^2}.
\]
这里上端速度为$\omega(a+b)$，轮心速度则为$\omega b$。把两者混同会使平移项多算一段。

再按平移与转动分开算。平移项为$B\omega b(a+b)$；转动项沿$y=-b$到$y=a$的直线参考路径为
\[B\omega\int_{-b}^{a}y\dd y
=B\omega\left[\frac{y^2}{2}\right]_{-b}^{a}
=\frac{B\omega}{2}(a^2-b^2).
\]
两者相加，仍为$B\omega(a+b)^2/2$，其中转动项为负。该直线只是一条便于积分的参考线，并不要求真实电流穿过轮毂；前面的梯度式已保证路径可更换。


\begin{figure}[H]\centering
\begin{tikzpicture}[>=Stealth,line cap=round,line join=round]
\begin{scope}[xshift=-2.4cm]\coordinate(A)at(0,-1.3);\coordinate(B)at(0,1.1);\coordinate(P)at(.9,1.1);
\draw[curve](A)--(B);\draw[guide](A)--(P);\draw[guide](B)--(P);\fill[gray!12](A)--(B)--(P)--cycle;
\node[below]at(A){$A$};\node[above]at(B){$B$};\node[right]at(P){$B'$};\draw[<->](-.4,-1.3)--(-.4,1.1)node[midway,left]{$a+b$};\node[below]at(0,-1.8){瞬心扫过三角形};\end{scope}
\begin{scope}[xshift=2.4cm]\draw[axis](0,-1.3)--(0,1.5)node[above]{$y$};\draw[axis](0,-1.3)--(2,-1.3)node[right]{$u_x$};\draw[curve](0,-1.3)--(1.6,1.1);\node[below left]at(0,-1.3){$-b$};\node[left]at(0,1.1){$a$};\draw[guide](0,1.1)--(1.6,1.1);\node[right]at(1.6,1.1){$\omega(a+b)$};\end{scope}
\end{tikzpicture}
\caption{端点速度和有效长度共同给出三角形扫掠面积。}\label{fig:i1-e9-s}
\end{figure}

图~\ref{fig:i1-e9-s}左图以瞬心$A$为轴。在$\dd t$内，上端扫过位移$u_B\dd t$，直线参考段扫过的面积为$(a+b)u_B\dd t/2$，乘磁场便得到同样的电动势。右图中$u_x=\omega(b+y)$的线性分布也给出同一个三角形面积。若用理想电刷接到外电路，应说明只有$A,B$两处导电，其他交叠处绝缘；几何电路换接时，直接积分比只看圆环面积是否变化更可靠。
\end{solution}

\begin{example}[内嵌无阻回路的圆柱能否一直滚下]\label{ex:i1-e10}
质量$m$、半径$r$、轴向长$h$的均匀实心圆柱放在倾角$0<\theta<\pi/2$的粗糙斜面上，静摩擦足以保持纯滚动。圆柱内固定一矩形线圈，跨越直径并沿轴向延伸，面积$2rh$，自感$L$、电阻为零。竖直匀强磁场大小为$B$；释放时线圈法向平行磁场，电流和柱速度均为零。求电流、运动方程及圆柱持续滚过所有势垒的磁场条件。 见图~\ref{fig:i1-e10-q}。
\end{example}

\begin{figure}[H]\centering
\begin{tikzpicture}[>=Stealth,line cap=round,line join=round]
\draw[curve](-2,-1.3)--(2.6,.25);\draw[curve](0,.1)circle(1.05);\draw[guide](-.95,-.35)--(.95,.55);\draw[->](0,.1)--(0,1.8)node[above]{$\vct B$};
\draw[->](1,.9)arc(45:-30:1.15)node[right]{$\psi$};\node[left]at(-1.1,.8){$m,r,h$};\node[below]at(0,-1.7){内嵌矩形回路：$S=2rh,\ L,\ R=0$};\node[right]at(1.5,-.65){$\theta$};
\end{tikzpicture}
\caption{回路随柱体转动，初始磁通最大；图为垂直圆柱轴的截面示意。}\label{fig:i1-e10-q}
\end{figure}
\begin{solution}
令$\psi$为图~\ref{fig:i1-e10-q}所示的滚转角，取下坡位移$x=r\psi$，矩形回路面积$S=2rh$。初始回路法向平行外磁场且电流为零，所以无阻回路的总磁链守恒给
\[LI+BS\cos\psi=BS,\qquad
\boxed{I=\frac{BS}{L}(1-\cos\psi)}.
\]
均匀外场对完整回路的合力为零，但有力矩。外磁力矩在滚转方向上的分量为$-IBS\sin\psi$，故它先阻碍滚下，在某些后续角度也可反向做功。

实心柱绕轴惯量$J=mr^2/2$。由纯滚动关系，平动与转动动能合为
\[T=\frac12m(r\dot\psi)^2+\frac12J\dot\psi^2
=\frac34mr^2\dot\psi^2.
\]
回路的自感能为$LI^2/2$。用电路式$LI\dot I=-I\dot\Phi_{\rm ext}$也可直接核对机械功与磁能的交换，静摩擦不做总功。由初始静止得
\[\frac34mr^2\dot\psi^2+
\frac{B^2S^2}{2L}(1-\cos\psi)^2-mgr\sin\theta\,\psi=0.
\]
记$K=B^2S^2/L$、$k=mgr\sin\theta/K$，以及$f(\psi)=(1-\cos\psi)^2/2$，则有效势为$U=K[f(\psi)-k\psi]$。图~\ref{fig:i1-e10-s}左图把越障问题化为：过原点的直线$k\psi$是否始终不低于$f(\psi)$。


\begin{figure}[H]\centering
\begin{tikzpicture}[>=Stealth,line cap=round,line join=round]
\begin{scope}[xshift=-2.8cm,x=.7cm,y=1cm]
\draw[axis](0,0)--(6.9,0)node[right]{$\psi$};\draw[axis](0,0)--(0,2.5)node[above]{$f(\psi)$};
\draw[curve,domain=0:6.283,samples=121]plot(\x,{(1-cos(57.29578*\x))^2/2});\draw[guide](0,0)--(3.5,2.3578)node[above]{$k_c\psi$};
\fill(2.7865,1.8772)circle(1.5pt);\node[below]at(2.7865,0){$\psi_c$};\end{scope}
\begin{scope}[xshift=2.8cm,x=.6cm,y=1cm]\draw[axis](0,0)--(6.8,0)node[right]{$\psi$};\draw[axis](0,-1.2)--(0,1.5)node[above]{$U/K$};
\draw[curve,domain=0:4.7,samples=101]plot(\x,{(1-cos(57.29578*\x))^2/2-.5*\x});\node[right]at(3.6,.4){小于临界：有势垒};\end{scope}
\end{tikzpicture}
\caption{过原点直线与周期函数的切触决定临界值，等号对应渐近越障极限。}\label{fig:i1-e10-s}
\end{figure}

因此临界斜率是
\[k_c=\max_{\psi>0}\frac{(1-\cos\psi)^2}{2\psi}.
\]
在$0<\psi<\pi$内，商法则给
\[\frac{\dd}{\dd\psi}\frac{(1-\cos\psi)^2}{2\psi}
=\frac{2\psi(1-\cos\psi)\sin\psi-(1-\cos\psi)^2}{2\psi^2}.\]
在非零驻点约去$1-\cos\psi>0$，得到
\[2\psi\sin\psi=1-\cos\psi,
\qquad\boxed{\tan\frac{\psi_c}{2}=2\psi_c}.
\]
该区间只有一个非零解。$\psi\ge\pi$时$f/\psi\le2/\pi$，低于这里的最大值，所以它也是全区间最大值。数值为
\[\psi_c\simeq2.78649815,\qquad k_c\simeq0.67366752.
\]
若$k<k_c$，柱从原点先加速，再在第一个满足$f(\psi)=k\psi$的点回转；若$k>k_c$，所有后续角度都可越过。等号时轨迹渐近接近势垒顶，不能在有限时间越过。于是从静止持续越过所有磁势垒要求
\[\boxed{B^2<\frac{mgL\sin\theta}{4rh^2k_c}}.
\]
图~\ref{fig:i1-e10-s}右图给出小于临界时的势垒，而非把临界角误认成任意参数下的停止角。

纯滚动还需静摩擦足够。取沿坡向上的摩擦力为$f_s$，在图~\ref{fig:i1-e10-q}的下坡、滚转正向分别列式：
\[mr\ddot\psi=mg\sin\theta-f_s,\qquad
\frac12mr^2\ddot\psi=f_sr-\tau_B,
\quad \tau_B=K(1-\cos\psi)\sin\psi.\]
第一式乘$r$再与第二式相加，消去$f_s$，即得所求运动方程
\[\boxed{\frac32mr^2\ddot\psi
=mgr\sin\theta-K(1-\cos\psi)\sin\psi},\qquad
\psi(0)=\dot\psi(0)=0.\]
于是$r\ddot\psi=\tfrac23g\sin\theta-2\tau_B/(3mr)$。代回$f_s=mg\sin\theta-mr\ddot\psi$，得到
\[f_s=\frac13mg\sin\theta+\frac{2\tau_B}{3r}.
\]
摩擦方向可能改变。题设“足够大”可用$|f_s|\le\mu_smg\cos\theta$逐点核对；它不是耗散力，不能另在能量式中减去摩擦功。
\end{solution}

\subsection{磁通与量子相位（选读）}
经典洛伦兹力由局域$E,B$决定，量子干涉还会涉及路径上的矢势积分。沿给定路径的相位变化含$\frac q\hbar\int\vct A\cdot\dd\vct l$；两条具有相同端点的路径相减，得到闭合环量。规范变换增加的两个端点项在路径差中抵消，因此可观测的是规范不变的相位差，而不是某一点的$A$值。

在粒子不能进入的细磁通管外，局域磁场可以为零，但绕管环量$\oint\vct A\cdot\dd\vct l=\Phi_B$仍不为零。若试图用$\chi\propto\theta$消去矢势，$\theta$绕一周会改变$2\pi$，这样的$\chi$不是该区域上的全域单值函数。势表示的全域性质因而进入干涉结果，这就是Aharonov--Bohm效应所讨论的构形。

\begin{example}[磁通管外的相位差]\label{ex:i1-e11}
电荷为$q$的粒子沿两条路径$C_1,C_2$从$S$到$D$，两路径围住一根不可进入的细磁通管，管内静态磁通为$\Phi$，粒子可达区域内$E=B=0$。已知路径相位中含$(q/\hbar)\int\vct A\cdot\dd\vct l$，求磁通引起的相位差和条纹复现的磁通周期。 见图~\ref{fig:i1-e11-q}。
\end{example}

\begin{figure}[H]\centering
\begin{tikzpicture}[>=Stealth,line cap=round,line join=round]
\fill[gray!25](0,0)circle(.35);\draw[curve](0,0)circle(.35);\node at(0,0){$\Phi$};\fill(-2,0)circle(1.5pt)node[left]{$S$};\fill(2,0)circle(1.5pt)node[right]{$D$};
\draw[curve,->](-2,0)..controls(-1.2,1.5)and(1.2,1.5)..(2,0);\draw[curve,->](-2,0)..controls(-1.2,-1.5)and(1.2,-1.5)..(2,0);
\node[above]at(0,1.2){$C_1$};\node[below]at(0,-1.2){$C_2$};\node[below]at(0,-1.8){管外$B=0$，粒子不进入阴影区};
\end{tikzpicture}
\caption{两路径组成的闭合曲线包围磁通管，路径本身不穿过磁场。}\label{fig:i1-e11-q}
\end{figure}
\begin{solution}
图~\ref{fig:i1-e11-q}的两条路径都在无磁场区域，单条路径上的矢势相位依赖规范；相位之差却为
\[\Delta\phi=\frac q\hbar\left(\int_{C_1}\vct A\cdot\dd\vct l-\int_{C_2}\vct A\cdot\dd\vct l\right)
=\frac q\hbar\oint_{C_1-C_2}\vct A\cdot\dd\vct l.
\]
闭合路径绕磁通管一周。按相应正向取$\Phi$，图~\ref{fig:i1-e11-s}给出的管外矢势$A_\theta=\Phi/(2\pi\rho)$满足局部旋度为零，但积分为$\Phi$，因此
\[\boxed{\Delta\phi=\frac{q\Phi}{\hbar}}.
\]

\begin{figure}[H]\centering
\begin{tikzpicture}[>=Stealth,line cap=round,line join=round]
\draw[curve](0,0)circle(1.25);\fill[gray!20](0,0)circle(.3);\foreach\a in{0,90,180,270}{\draw[->](\a:1.25)arc(\a:\a+45:1.25);}
\node[right]at(1.6,.5){$A_\theta=\Phi/(2\pi\rho)$};\node[right]at(1.6,-.1){$\oint A_\theta\rho\dd\theta=\Phi$};\node[right]at(1.6,-.7){$\Delta\phi=q\Phi/\hbar$};
\end{tikzpicture}
\caption{局部无旋与全域非零环量可以同时成立。}\label{fig:i1-e11-s}
\end{figure}

规范变换下$\vct A'=\vct A+\nabla\chi$。两条开路径上的积分各增加$\chi(D)-\chi(S)$，所以
\[
\left[\int_{C_1}\nabla\chi\cdot\dd\vct l
-\int_{C_2}\nabla\chi\cdot\dd\vct l\right]
=[\chi(D)-\chi(S)]-[\chi(D)-\chi(S)]=0.
\]
这说明相位差与单值规范函数的选择无关。沿图中的外圆积分时，
$\oint A_\theta r\dd\theta=\frac{\Phi}{2\pi}[\theta]_0^{2\pi}=\Phi$，不能因为各处$B=0$就把跨过不可达磁通管的曲面也当作无场面。

干涉条纹在相位增加$2\pi$时复现，所以磁通周期为
\[\boxed{\Delta\Phi=\frac{2\pi\hbar}{|q|}=\frac h{|q|}}.
\]
电子的$q=-e$改变条纹移动方向，不改变周期的正值。粒子在管外不受经典磁力，磁通仍可通过量子相位影响干涉。这里使用了量子相位规律及磁通区不可进入的条件；可观测量是两路径的相位差，矢势的局部数值依赖规范。若两条路径围不住磁通管，闭合环量为零，此处讨论的磁通相位差也消失。
\end{solution}

\section{感应电压的测量路径与滑动接触}
\subsection{非保守电场中的两端电压}
静电场中，两点间电压只由端点决定。存在变化磁通时，$\oint\vct E\cdot\dd\vct l\ne0$，实际测量还必须给出引线怎样连接。取两条从$A$到$B$的静止路径$\gamma_1,\gamma_2$，定义沿各路径的有向电场积分
\[
 U_1=\int_{\gamma_1,A\to B}\vct E\cdot\dd\vct l,\qquad
 U_2=\int_{\gamma_2,A\to B}\vct E\cdot\dd\vct l.
\]
闭合曲线按$\gamma_1$正向、$\gamma_2$反向组成，其磁通正向由右手规则确定。于是
\[
 U_1-U_2=-\frac{\dd\Phi_{\gamma_1-\gamma_2}}{\dd t}.
\]
理想高内阻电压表通过引线接入时，对原回路的扰动可忽略。在静止、细导线及准静态条件下，读数可用沿实际引线和表内路径的积分表示。改变引线包围的磁通，会改变读数。将两根引线并行靠近或绞合，可减小附加回路面积；仅仅说“表接在$A,B$两端”还不足以确定答案。

\begin{example}[同一对接点的两个电压表读数]
一静止圆形回路分成上、下两条半圆支路，接点分别为左侧$A$、右侧$B$。上支路电阻为$2R$，下支路电阻为$R$，自感可忽略。穿过圆面的向纸外磁通$\Phi(t)$均匀增加，$\dot\Phi=\mathcal E_0>0$。两个高内阻电压表均以$A$端为正端、$B$端为负端；第一组引线紧贴上支路，第二组紧贴下支路，分别与对应支路包围可忽略的小面积。求两表读数。
\ExpansionDiagram{fig:dep-i1-meter-q}{两组引线分别沿上、下支路连接同一对接点}{
\draw[thick](-2,0)arc(180:115:2);\draw[thick](2,0)arc(0:65:2);
\draw[thick](-.845,1.813)--(-.6,1.65)--(-.3,2.05)--(0,1.65)--(.3,2.05)--(.6,1.65)--(.845,1.813);
\draw[thick](-2,0)arc(180:245:2);\draw[thick](2,0)arc(360:295:2);
\draw[thick](-.845,-1.813)--(-.6,-1.65)--(-.3,-2.05)--(0,-1.65)--(.3,-2.05)--(.6,-1.65)--(.845,-1.813);
\fill(-2,0)circle(2pt)node[left=12pt]{$A$};\fill(2,0)circle(2pt)node[right=12pt]{$B$};
\node[fill=white]at(0,1.45){$2R$};\node[fill=white]at(0,-1.45){$R$};
\draw[dashed,->](-2.3,.1)arc(177.5:2.5:2.3);\draw[dashed,->](-2.3,-.1)arc(182.5:357.5:2.3);
\node[right]at(2.3,1.3){$\gamma_1$};\node[right]at(2.3,-1.3){$\gamma_2$};
\node at(0,0){$\odot\ \dot\Phi>0$};}
\end{example}
\begin{solution}
图\ref{fig:dep-i1-meter-q}中的向外磁通增加，感应电流为顺时针。在上半圆，电流由$A$流向$B$；在下半圆，由$B$流向$A$。总电阻$3R$，所以
\[
 I=\frac{\mathcal E_0}{3R}.
\]
第一组引线与上支路之间磁通变化可忽略，沿两路径的电场积分相同。上支路从$A$到$B$与电流同向，因而
\[
 U_1=\int_{A\to B,\text{上}}\vct E\cdot\dd\vct l=2RI
 =\frac23\mathcal E_0.
\]
第二组引线紧贴下支路，从$A$到$B$的积分方向逆着电流，故
\[
 U_2=\int_{A\to B,\text{下}}\vct E\cdot\dd\vct l=-RI
 =-\frac13\mathcal E_0.
\]
\ExpansionDiagram{fig:dep-i1-meter-a}{顺时针电流与两条积分路径的符号}{
\draw[thick](-2,0)arc(180:0:2);\draw[thick](-2,0)arc(180:360:2);
\draw[->,thick](-1.414,1.414)arc(135:65:2)node[midway,above right]{$I$};
\draw[->,thick](1.414,-1.414)arc(315:245:2)node[midway,below right]{$I$};
\node[left]at(-2,0){$A$};\node[right]at(2,0){$B$};
\draw[->](-1,.6)--(1,.6)node[midway,above]{$U_1=+2RI$};
\draw[->](-1,-.6)--(1,-.6)node[midway,below]{$U_2=-RI$};}
按图\ref{fig:dep-i1-meter-a}中的上路$A\to B$再下路$B\to A$构成顺时针闭合曲线，其正法向指向纸内，磁通为$-\Phi$。于是
$U_1-U_2=-\dd(-\Phi)/\dd t=\mathcal E_0$，与两个读数相符。电压表理想内阻趋于无穷不消除路径差异，只使它对原回路电流的影响趋于零。
\end{solution}

\subsection{滑动接触回路怎样使用磁通法则}
运动回路的磁通导数公式以连续随物质运动的闭合曲线为对象。电刷、滑轨等滑动接触会不断更换路径上的物质点；处理这类回路，应按各段的实际导体速度分别计算
\[
 \mathcal E=\oint(\vct E+\vct v_{\rm wire}\times\vct B)\cdot\dd\vct l,
\]
再说明各段导体速度及接触处连接方式。不能把某个几何上不动的回路面积直接代入$-\dd\Phi/\dd t$，同时漏掉运动导体段的贡献。

\begin{example}[匀强静磁场中的单极发电圆盘]
半径$a$的导电圆盘在垂直盘面向外的匀强磁场$B$中，以角速度$\omega>0$逆时针转动。静止电刷分别接触圆心$A$和边缘$B_1$，外电阻为$R$，圆盘及接触的总内阻为$r$。忽略接触的局部效应。求开路电压、闭路电流、阻力矩和维持转动所需的功率。
\ExpansionDiagram{fig:dep-i1-disk-q}{圆心和盘边电刷接入静止外回路}{
\draw[thick](0,0)circle(1.5);\fill(0,0)circle(2pt)node[below left]{$A$};\fill(1.5,0)circle(2pt)node[above right]{$B_1$};
\draw[thick](0,0)--(0,-2.2)--(2.8,-2.2);\draw[thick](1.5,0)--(2.8,0)--(2.8,-.7);
\draw[thick](2.8,-.7)--(2.65,-.9)--(2.95,-1.1)--(2.65,-1.3)--(2.95,-1.5)--(2.8,-1.7)--(2.8,-2.2);
\node[right]at(3,-1.2){$R$};\node at(-.7,.3){$\odot B$};
\draw[->](1.1,.8)arc(36:125:1.36);\node[above]at(0,1.7){$\omega$};}
\end{example}
\begin{solution}
图\ref{fig:dep-i1-disk-q}中盘上距圆心$s$的质点速度为$\omega s\,\vct e_\varphi$。由于$\vct e_\varphi\times\vct e_z=\vct e_s$，磁力将正电荷推向外缘。沿盘上半径从$A$到$B_1$，运动电动势为
\[
 \mathcal E_{\rm d}=\int_0^a\omega sB\,\dd s
 =B\omega[s^2/2]_0^a
 =\frac12B\omega a^2.
\]
开路时，盘内静电场调整为$E_s=-B\omega s$，所以
\[
 \boxed{\varphi_{B_1}-\varphi_A=\frac12B\omega a^2}.
\]
闭路后电流在盘中由内向外，经过外电阻回到圆心，其大小
$I=B\omega a^2/[2(R+r)]$。在任意半径$s$，穿过同心圆的总径向电流都是$I$。若盘厚为$h$，则$J_s=I/(2\pi sh)$；薄环体积为$2\pi sh\,\dd s$。由$\vct J\times\vct B$，该环受到的切向力为$-IB\,\dd s$，力矩为
\[
 \dd\tau_z=-IBs\,\dd s,\qquad
 \tau_z=-IB\int_0^as\dd s=-IB[s^2/2]_0^a,\qquad
 \boxed{\tau_z=-\frac12IBa^2}.
\]
\ExpansionDiagram{fig:dep-i1-disk-a}{径向电流、切向磁力及机械能转为电能的方向}{
\draw[thick](0,0)circle(1.7);\draw[dashed](0,0)circle(1);
\draw[->,thick](.2,0)--(1.55,0)node[midway,above]{$I$};
\fill(1,0)circle(2pt);\draw[->,thick](1,0)--(1,-1.15)node[right]{$\dd F_\varphi$};
\draw[<->](0,-.25)--(1,-.25)node[midway,below]{$s$};
\node at(-.7,.4){$\odot B$};\node[below]at(0,-2){$\tau_z<0,\quad \omega>0$};}
图\ref{fig:dep-i1-disk-a}中的磁力矩反向于圆盘转动。保持角速度恒定，外力矩应为$IBa^2/2$，提供的功率为
\[
 P_{\rm mech}=\frac12IBa^2\omega=I\mathcal E_{\rm d}
 =I^2(R+r).
\]
其中$I^2R$供给外负载，$I^2r$转为内阻热。整个几何电路处在静止磁场中，几何围成的面积也可不变，电流仍然存在；圆盘段的速度场和滑动电刷使这条路径不是物质回路。
\end{solution}

\section{习题与解答}

\begin{exercise}[磁场减弱而总磁通不变]\label{ex:i1-h1}
一个半径可变的材料圆环位于$xy$平面，外磁场为$B_0\ee^{-t/\tau}\vct e_z$。初始半径$a_0$，要求总电动势始终为零。求$a(t)$，并分别写出感生与动生贡献。 见图~\ref{fig:i1-h1-q}。
\end{exercise}

\begin{figure}[H]\centering
\begin{tikzpicture}[>=Stealth,line cap=round,line join=round]
\draw[curve](0,0)circle(1.3);\draw[guide](0,0)circle(1.65);\draw[->](1.3,0)--(1.65,0);\node at(0,0){$\odot\ B_0\ee^{-t/\tau}$};\node[above]at(0,1.9){$a(t)$};
\end{tikzpicture}
\caption{虚线表示稍晚时刻的材料圆环。}\label{fig:i1-h1-q}
\end{figure}
\begin{solution}
沿图~\ref{fig:i1-h1-q}的逆时针方向积分，磁通为$\Phi=\pi a^2B_0\ee^{-t/\tau}$。总电动势为零要求其全时间导数为零，而不是要求磁场和面积分别不变。因此
\[\frac{\dd\Phi}{\dd t}=\pi B_0\ee^{-t/\tau}
\left(2a\dot a-\frac{a^2}{\tau}\right)=0,\qquad
\frac{\dot a}{a}=\frac1{2\tau}.
\]
从初始半径$a_0$与时刻$0$分别积分，
\[\int_{a_0}^{a(t)}\frac{\dd s}s=\int_0^t\frac{\dd t'}{2\tau},\qquad
[\ln s]_{a_0}^{a(t)}=\frac{t}{2\tau},\qquad
\ln\frac{a(t)}{a_0}=\frac{t}{2\tau}.\]
取指数得
\[\boxed{a(t)=a_0\ee^{t/(2\tau)}}.
\]

\begin{figure}[H]\centering
\begin{tikzpicture}[>=Stealth,line cap=round,line join=round]
\draw[axis](-.1,0)--(4.5,0)node[right]{$t/\tau$};\draw[axis](0,-.1)--(0,2.5)node[above]{$a/a_0$};
\draw[curve,domain=0:2.1,samples=71]plot(\x,{exp(\x/2)});
\node[above left]at(1.6,2.25){$\ee^{t/(2\tau)}$};
\end{tikzpicture}
\caption{半径的增长时标是磁场衰减时标的两倍。}\label{fig:i1-h1-s}
\end{figure}

图~\ref{fig:i1-h1-s}画出的扩大速度越来越大。场的减弱与面积的增加严格补偿，使$\pi a^2B=\pi a_0^2B_0$保持常量。

分别计算两种贡献也得到相同结论：感生项$\mathcal E_{\rm tr}=+\pi a^2B/\tau$，动生项$\mathcal E_{\rm mot}=-2\pi aB\dot a=-\pi a^2B/\tau$。因此忽略自感且初始无电流时，电阻环中始终没有环流。环上局部仍可能有非零电场与表面电荷分布；总电动势为零只约束沿整个回路的积分。若外磁场只在有限磁区内满足给定规律，解还要求圆环在所讨论时段内没有越出该区域。
\end{solution}

\begin{exercise}[环形磁区与空心区域]\label{ex:i1-h2}
轴对称磁场仅在$a<r<b$的圆环区域随时间变化，$\partial_tB_z=-\alpha$；其他区域$\partial_tB_z=0$。求全部三个区域的感生电场，并讨论两个边界与$a\to0$极限。 见图~\ref{fig:i1-h2-q}。
\end{exercise}

\begin{figure}[H]\centering
\begin{tikzpicture}[>=Stealth,line cap=round,line join=round]
\draw[curve](0,0)circle(1.6);\draw[curve](0,0)circle(.75);\fill[even odd rule,gray!12](0,0)circle(1.6)(0,0)circle(.75);\draw[axis](-2,0)--(2.1,0)node[right]{$x$};\node at(0,.25){$O$};\node at(1.1,.75){$-\dot B_z=\alpha$};\node[below]at(.75,0){$a$};\node[below]at(1.6,0){$b$};
\end{tikzpicture}
\caption{阴影为变化磁场的截面，内孔不含该场。}\label{fig:i1-h2-q}
\end{figure}
\begin{solution}
图~\ref{fig:i1-h2-q}中的变化磁场只存在于环形区域。对半径$r$的逆时针圆周用Faraday定律，左端均为$2\pi rE_\theta$，但右端的有效面积分三段。

当$r<a$时，没有包围变化磁通，$E_\theta=0$。当$a<r<b$时，包围面积为$\pi(r^2-a^2)$；当$r>b$时，该面积固定为$\pi(b^2-a^2)$。所以
\[\boxed{E_\theta(r)=\begin{cases}
0,&r<a,\\
\dfrac\alpha2\left(r-\dfrac{a^2}{r}\right),&a<r<b,\\
\dfrac{\alpha(b^2-a^2)}{2r},&r>b.
\end{cases}}
\]

\begin{figure}[H]\centering
\begin{tikzpicture}[>=Stealth,line cap=round,line join=round]
\draw[axis](-.1,0)--(4.5,0)node[right]{$r$};\draw[axis](0,-.1)--(0,2.5)node[above]{$E_\theta$};
\draw[curve](0,0)--(1,0);
\draw[curve,domain=1:2,samples=60]plot(\x,{.9*(\x-1/\x)});
\draw[curve,domain=2:4.1,samples=60]plot(\x,{2.7/\x});
\node[below]at(1,0){$a$};\node[below]at(2,0){$b$};\draw[guide](2,0)--(2,1.35);

\end{tikzpicture}
\caption{积分面积分成三段，电场在两个边界连续。}\label{fig:i1-h2-s}
\end{figure}

图~\ref{fig:i1-h2-s}中两处连接都连续：在$a$处为零，在$b$处为$\alpha(b^2-a^2)/(2b)$。圆周切向场不需要发生跳跃。

区内旋度也可分段复算。在环形区域，
\[\frac1r\frac{\dd(rE_\theta)}{\dd r}
=\frac\alpha{2r}\frac{\dd(r^2-a^2)}{\dd r}
=\frac\alpha{2r}(2r)=\alpha.\]
其他两区的$rE_\theta$都是常数，径向导数为零。内孔不含变化磁通，所以中心附近没有感生场；外区虽没有局域变化磁场，却包围了整个变化环带，故电场不为零。令$a\to0$即可恢复实心圆磁区的内外解，这个极限也核对了分母与因子$1/2$。
\end{solution}

\begin{exercise}[磁场变化率随半径增大]\label{ex:i1-h3}
在$r<R$内，$\partial_tB_z=-\alpha r^2/R^2$，区外为零。求感生电场，比较它与均匀变化圆磁区在边界处的场，并用局部方程核对。 见图~\ref{fig:i1-h3-q}。
\end{exercise}

\begin{figure}[H]\centering
\begin{tikzpicture}[>=Stealth,line cap=round,line join=round]
\fill[gray!8](0,0)circle(1.45);\draw[curve](0,0)circle(1.45);
\draw[axis](-1.9,0)--(2.1,0)node[right]{$x$};\draw[axis](0,-1.8)--(0,1.9)node[above]{$y$};
\node[below left]at(0,0){$O$};\draw[->](0,0)--(40:1.45)node[midway,above left]{$R$};
\node at(-.65,.65){$\dot B_z$};\fill(1,0)circle(1.5pt)node[below right]{$P$};\node[below]at(0,-2.1){$-\dot B_z=\alpha r^2/R^2\quad(r<R)$};
\end{tikzpicture}
\caption{边缘变化率与均匀模型相同，盘内平均值不同。}\label{fig:i1-h3-q}
\end{figure}
\begin{solution}
图~\ref{fig:i1-h3-q}仍有轴对称性，但磁通变化率在圆内不均匀。对$r<R$的圆盘，应逐个环带积分：
\[2\pi rE_\theta=\int_0^r\frac{\alpha s^2}{R^2}\,2\pi s\dd s
=\frac{2\pi\alpha}{R^2}\left[\frac{s^4}4\right]_0^r
=\frac{\alpha\pi r^4}{2R^2}.
\]
从而$E_\theta=\alpha r^3/(4R^2)$。在$r>R$时把积分上限固定为$R$，得到$2\pi rE_\theta=\alpha\pi R^2/2$。全解为
\[\boxed{E_\theta=\begin{cases}
\alpha r^3/(4R^2),&r<R,\\
\alpha R^2/(4r),&r>R.
\end{cases}}
\]

\begin{figure}[H]\centering
\begin{tikzpicture}[>=Stealth,line cap=round,line join=round]
\draw[axis](-.1,0)--(4.5,0)node[right]{$r/R$};\draw[axis](0,-.1)--(0,2.5)node[above]{$4E_\theta/(\alpha R)$};
\draw[curve,domain=0:1,samples=51]plot(\x,{(\x)^3});\draw[curve,domain=1:4,samples=61]plot(\x,{1/\x});
\draw[guide](1,0)--(1,1);\node[below]at(1,0){$1$};
\end{tikzpicture}
\caption{三次幂内解在边界接上反比外解。}\label{fig:i1-h3-s}
\end{figure}

图~\ref{fig:i1-h3-s}表明内区先按三次幂增长，再接外区的反比曲线。边界值为$\alpha R/4$，只有同样边界变化率的均匀圆磁区的一半。

差别来自面积平均：圆盘内$r^2/R^2$的面积平均为$1/2$，以边缘变化率代替面内平均值会多算一倍。将内解代入旋度式，
\[\frac1r\frac{\partial(rE_\theta)}{\partial r}
=\frac\alpha{4R^2r}\frac{\dd r^4}{\dd r}
=\frac\alpha{4R^2r}\,4r^3=\frac{\alpha r^2}{R^2},\]
恢复给定的局部源；外解的$rE_\theta$为常量，旋度为零。一般轴对称分布也应先求$\int_0^r[-\dot B_z(s)]s\dd s$，再除以$r$。
\end{solution}

\begin{exercise}[从圆磁区的矢势恢复感生场]\label{ex:i1-h4}
半径$R$圆柱内的匀强磁场为$B(t)\vct e_z$，区外磁场为零。选仅有环向分量的磁矢势，要求轴上正则。求内外$A_\theta$，再在无附加电荷场的约定下恢复$E_\theta$。 见图~\ref{fig:i1-h4-q}。
\end{exercise}

\begin{figure}[H]\centering
\begin{tikzpicture}[>=Stealth,line cap=round,line join=round]
\draw[curve](0,0)circle(1.5);\draw[axis](-1.9,0)--(2,0);\draw[->](0,0)--(45:1.5)node[midway,above]{$R$};\node at(-.65,.6){$\odot\ B(t)$};\draw[->](1.2,0)arc(0:70:1.2)node[above]{$A_\theta$};
\end{tikzpicture}
\caption{环向矢势的环量必须给出圆柱内的磁通。}\label{fig:i1-h4-q}
\end{figure}
\begin{solution}
对图~\ref{fig:i1-h4-q}选$\vct A=A_\theta(r,t)\vct e_\theta$，并要求轴上无奇点。由柱坐标旋度，
\[B_z=\frac1r\frac{\partial}{\partial r}(rA_\theta).
\]
固定时刻后，圆内满足$\partial_r(rA_\theta)=rB(t)$，故
\[rA_\theta(r,t)-\lim_{s\to0}sA_\theta(s,t)
=B(t)\int_0^rs\dd s=B(t)[s^2/2]_0^r.\]
轴上正则使左端极限为零，于是$A_\theta=Br/2$。圆外$\partial_r(rA_\theta)=0$，通解为$A_\theta=C_o(t)/r$；由$r=R$处连续得$C_o=BR^2/2$。


\begin{figure}[H]\centering
\begin{tikzpicture}[>=Stealth,line cap=round,line join=round]
\draw[axis](-.1,0)--(4.5,0)node[right]{$r/R$};\draw[axis](0,-.1)--(0,2.5)node[above]{$2A_\theta/(BR)$};
\draw[curve](0,0)--(1,1);\draw[curve,domain=1:4,samples=61]plot(\x,{1/\x});
\node[below]at(1,0){$1$};\node[right]at(2.1,1.5){$E_\theta=-\partial_tA_\theta$};
\end{tikzpicture}
\caption{矢势的内外分段与磁通相容，其时间导数给出感生场。}\label{fig:i1-h4-s}
\end{figure}

因此图~\ref{fig:i1-h4-s}的两段为
\[\boxed{A_\theta=\begin{cases}Br/2,&r<R,\\BR^2/(2r),&r>R.\end{cases}}
\]
在这个对称、无附加电荷场的规范下可取$\varphi=0$，于是$E_\theta=-\partial_tA_\theta$，立即得到圆内$-\dot Br/2$、圆外$-\dot BR^2/(2r)$。

外区$B=0$并不允许把$A_\theta$也随意取零。沿外圆的矢势环量为$2\pi rA_\theta=\pi R^2B$，等于管内磁通；若改成零便改变了这一全域量。只有局部、可缩的小邻域内，才可把这样的无旋矢势写成梯度。计算感生电场时，还必须保留与规范同时变换的电势。
\end{solution}

\begin{exercise}[时间依赖规范不能只改变矢势]\label{ex:i1-h5}
匀强磁场$B(t)\vct e_z$采用对称规范$\vct A=\tfrac12B(-y,x,0)$、$\varphi=0$。改用$\vct A'=(-By,0,0)$，求规范函数和新电势，并逐分量说明总电场不变。 见图~\ref{fig:i1-h5-q}。
\end{exercise}

\begin{figure}[H]\centering
\begin{tikzpicture}[>=Stealth,line cap=round,line join=round]
\draw[axis](-1.8,0)--(2,0)node[right]{$x$};\draw[axis](0,-1.4)--(0,1.8)node[above]{$y$};\fill(1,1)circle(1.5pt)node[above right]{$P$};\node at(-1,.8){$\odot\ B(t)$};\draw[->](1,1)--(1.5,.5);\node[below]at(0,-1.8){两种势表示同一个物理电场};
\end{tikzpicture}
\caption{取一般点P，比较电场矢量而非某个势的数值。}\label{fig:i1-h5-q}
\end{figure}
\begin{solution}
图~\ref{fig:i1-h5-q}中的对称规范为$\vct A=\tfrac12B(-y,x,0)$、$\varphi=0$。要变到$\vct A'=(-By,0,0)$，两矢势之差为$(-By/2,-Bx/2,0)$，所以可取
\[\chi=-\frac12B(t)xy,\qquad
\varphi'=\varphi-\partial_t\chi=\frac12\dot Bxy.
\]
在原规范中，$\vct E=-\partial_t\vct A=(\dot By/2,-\dot Bx/2,0)$。


\begin{figure}[H]\centering
\begin{tikzpicture}[>=Stealth,line cap=round,line join=round]
\draw[axis](-.2,0)--(3.8,0)node[right]{$E_x$};\draw[axis](0,-1.6)--(0,1.3)node[above]{$E_y$};
\draw[->,thick](0,0)--(2,-1)node[right]{$-\partial_t\vct A-\nabla\varphi$};\draw[->](0,0)--(3,0)node[above]{$-\partial_t\vct A'$};\draw[->](3,0)--(2,-1)node[midway,right]{$-\nabla\varphi'$};
\end{tikzpicture}
\caption{新规范中的梯度场补偿矢势时间导数的变化。}\label{fig:i1-h5-s}
\end{figure}

图~\ref{fig:i1-h5-s}将新规范中的两项分别画出。逐项计算，
\[-\partial_t\vct A'=(\dot By,0,0),\qquad
-\nabla\varphi'=(-\dot By/2,-\dot Bx/2,0).
\]
相加仍为
\[\boxed{\vct E=(\dot By/2,-\dot Bx/2,0)}.
\]
规范变换重新分配了电场的两项表示，新电势恰好补偿矢势改变的部分，物理电荷分布保持不变。总场的散度仍为零，旋度为$-\dot B\vct e_z$。若在新规范中仍令$\varphi'=0$，所得$(\dot By,0,0)$虽有同一旋度，却改变了总电场。Faraday方程给定旋度后，电场的梯度部分还须由其他条件确定。
\end{solution}

\begin{exercise}[非均匀静磁场中的有效势]\label{ex:i1-h6}
在平面内，$\vct B=B_0(1+y^2/a^2)\vct e_z$，$B_0>0$。质量$m$、电荷$q\ne0$的粒子无外力运动。用循环坐标将轨迹约化为一个$y$方向有效势问题，写出转向点条件。 见图~\ref{fig:i1-h6-q}。
\end{exercise}

\begin{figure}[H]\centering
\begin{tikzpicture}[>=Stealth,line cap=round,line join=round]
\draw[axis](-1.8,0)--(2,0)node[right]{$x$};\draw[axis](0,-1.6)--(0,1.7)node[above]{$y$};\foreach\y in{-1,-.5,0,.5,1}{\node at(-1.1,\y){$\odot$};}\node[right]at(.35,1.25){$B_z=B_0(1+y^2/a^2)$};\draw[->](.3,.3)--(1.2,.7)node[right]{$\vct v_0$};
\end{tikzpicture}
\caption{静磁场保持速率，不保持回旋半径。}\label{fig:i1-h6-q}
\end{figure}
\begin{solution}
图~\ref{fig:i1-h6-q}中的磁场只依赖$y$，可选$A_x=-F(y)$、$A_y=0$，其中
\[F(y)=\int_0^yB_z(s)\dd s
=B_0\left[s+\frac{s^3}{3a^2}\right]_0^y
=B_0\left(y+\frac{y^3}{3a^2}\right).
\]
拉格朗日量不显含$x$，因此$P_x=mv_x-qF(y)=p$为常量，且磁场不做功，平面总动能$E_\perp$为常量。消去$v_x$，
\[\frac12m\dot y^2+\frac{[p+qF(y)]^2}{2m}=E_\perp.
\]

\begin{figure}[H]\centering
\begin{tikzpicture}[>=Stealth,line cap=round,line join=round]
\draw[axis](-.1,0)--(4.5,0)node[right]{$y$};\draw[axis](0,-.1)--(0,2.5)node[above]{$U_{\rm eff}$};
\draw[curve,domain=-1.2:1.2,samples=91]plot({1.8+\x},{.7*(\x+(\x)^3/3)^2});\draw[guide](.3,1)--(3.4,1);
\node[right]at(3.4,1){$E_\perp$};\node[below]at(.8,0){$y_-$};\node[below]at(2.8,0){$y_+$};
\end{tikzpicture}
\caption{图示p=0的有效势，水平能量线确定两个转向点。}\label{fig:i1-h6-s}
\end{figure}

图~\ref{fig:i1-h6-s}中的一维有效势为
\[\boxed{U_{\rm eff}(y)=\frac{[p+qB_0(y+y^3/3a^2)]^2}{2m}}.
\]
转向点满足$p+qF(y)=\pm\sqrt{2mE_\perp}$。因为$F'(y)=B_z(y)>0$且$F$的值域为全体实数，当$E_\perp>0$时，两种符号各对应一个实根，两个根之间是允许的$y$区间。若$E_\perp=0$，两根合并为$F(y_0)=-p/q$，平面速度为零，粒子保持静止。图中取$p=0$以突出对称形状。一般$p$使势阱最低点移至$y_0$；$F$的非线性还会改变势阱形状。

对有效势求导可核对约化后的动力学：
\[U_{\rm eff}'(y)=\frac{2[p+qF(y)]qF'(y)}{2m}
=qB_z(y)v_x,\qquad m\ddot y=-qB_z(y)v_x.\]
这与图~\ref{fig:i1-h6-q}中洛伦兹力的$y$分量相同。在最低点再求一次导数，乘积法则给
\[U_{\rm eff}''(y_0)
=\left.\frac{q^2[F'(y)]^2+q[p+qF(y)]F''(y)}m\right|_{y_0}
=\frac{q^2B_z(y_0)^2}{m}.\]
因此最低点附近的曲率也取决于$p$所确定的$y_0$，不能将整族曲线理解为同一曲线的平移。

这里不能仍把轨道当作固定半径圆：磁场的大小随位置改变，机械速度方向的角速率也改变。所得有效势并非真实静电势，而是将守恒的正则动量代回机械能后形成的单坐标描述。允许区间内再由$v_x=[p+qF(y)]/m$与$\dot y$的正负支积分，才能得到具体轨迹。
\end{solution}

\begin{exercise}[由中心电场反求磁区长宽比]\label{ex:i1-h7}
无附加电荷场的矩形磁区内$\partial_tB_z=-\alpha$。在中心附近测得$E_x=-ay+O(r^3)$、$E_y=bx+O(r^3)$，且$a=3b$。求矩形沿$x$方向全长与沿$y$方向全宽之比，并说明测量范围的要求。 见图~\ref{fig:i1-h7-q}。
\end{exercise}

\begin{figure}[H]\centering
\begin{tikzpicture}[>=Stealth,line cap=round,line join=round]
\draw[curve](-2,-.7)rectangle(2,.7);\draw[axis](-2.3,0)--(2.4,0)node[right]{$x$};\draw[axis](0,-1)--(0,1.2)node[above]{$y$};\fill(.3,.3)circle(1.4pt)node[above right]{$P$};\node[below]at(0,-1.3){中心线性场：$E_x=-ay,\ E_y=bx$};
\end{tikzpicture}
\caption{用中心极限的系数比反演几何形状。}\label{fig:i1-h7-q}
\end{figure}
\begin{solution}
图~\ref{fig:i1-h7-q}给的是中心场的线性系数。局部Faraday方程要求$a+b=\alpha$，而$a=3b$，所以$a=3\alpha/4$、$b=\alpha/4$。

设矩形的两半边长为$D/2,H/2$。例~\ref{ex:i1-e5}已经由磁区面积核积分得到中心系数，直接用其中的关系
\[a=\frac{2\alpha}{\pi}\arctan\frac DH,\qquad
b=\frac{2\alpha}{\pi}\arctan\frac HD.\]
这里反演的是同一中心极限，无须再作一次面积核积分。

将$a/\alpha=3/4$代入，得到
\[\arctan\frac DH=\frac{3\pi}{8},\qquad
\boxed{\frac DH=\tan\frac{3\pi}{8}=1+\sqrt2}.
\]

\begin{figure}[H]\centering
\begin{tikzpicture}[>=Stealth,line cap=round,line join=round]
\draw[axis](-.1,0)--(4.5,0)node[right]{$D/H$};\draw[axis](0,-.1)--(0,2.5)node[above]{$a/\alpha$};
\draw[curve,domain=.05:4.2,samples=91]plot(\x,{2*atan(\x)/180});
\draw[guide](2.414,0)--(2.414,.75)--(0,.75);\node[below]at(2.414,0){$1+\sqrt2$};\node[left]at(0,.75){$3/4$};
\end{tikzpicture}
\caption{纵坐标采用弧度定义的反正切函数，交点给出长宽比。}\label{fig:i1-h7-s}
\end{figure}

图~\ref{fig:i1-h7-s}的水平线在$3/4$处与系数曲线相交，由此确定长宽比。磁场变化率的绝对大小则由$\alpha=a+b$给出。

测量点必须足够接近中心，使三次及更高项可忽略。若在同一个有限离轴点用$-E_x/y$和$E_y/x$当作中心系数，得到的比值会受高阶项污染。实际可在几个递减距离上作线性拟合，再外推到中心；不能把上述局部反演公式用于整个矩形边界。另有附加静电场时，也须先分离该场，否则同一个磁区形状未必给出这里的系数。
\end{solution}

\begin{exercise}[干涉路径的绕数]\label{ex:i1-h8}
电子沿两条无磁场路径从同一点到同一终点，绕不可进入的磁通管的有向次数分别为$2$和$-1$。求磁通引起的相位差与条纹复现周期，并说明周期为什么不能单独决定载流子电荷。 见图~\ref{fig:i1-h8-q}。
\end{exercise}

\begin{figure}[H]\centering
\begin{tikzpicture}[>=Stealth,line cap=round,line join=round]
\fill[gray!20](0,0)circle(.3);\node at(0,0){$\Phi$};\draw[curve,->](0,1.1)arc(90:420:1.1);\draw[curve,->](0,1.4)arc(90:420:1.4);\draw[guide,->](0,-1.7)arc(-90:-420:1.7);\node[right]at(2,.7){$n_1=2$};\node[right]at(2,0){$n_2=-1$};
\end{tikzpicture}
\caption{实线与虚线只示意绕数，端点相同。}\label{fig:i1-h8-q}
\end{figure}
\begin{solution}
图~\ref{fig:i1-h8-q}中的两条路径具有相同端点，绕管次数分别为$2$和$-1$。相减得到的闭合曲线不是只绕一周，而是净绕$\Delta n=3$周，所以其矢势环量为$3\Phi$。

电子$q=-e$的相位差为$\Delta\phi=-3e\Phi/\hbar$。若只关心条纹复现的磁通增量，应取相位变化的绝对值为$2\pi$，得到
\[\boxed{\Delta\Phi=\frac{2\pi\hbar}{3e}=\frac h{3e}}.
\]

\begin{figure}[H]\centering
\begin{tikzpicture}[>=Stealth,line cap=round,line join=round]
\draw[axis](-.1,0)--(4.5,0)node[right]{$\Phi$};\draw[axis](0,-.1)--(0,2.5)node[above]{$\Delta\phi$};
\draw[curve](0,0)--(3.6,2.1);\draw[guide](0,1.2)--(2.057,1.2)--(2.057,0);
\node[left]at(0,1.2){$2\pi$};\node[below]at(2.057,0){$h/(3e)$};
\end{tikzpicture}
\caption{净绕数为三，磁通周期缩为单绕数构形的三分之一。}\label{fig:i1-h8-s}
\end{figure}

图~\ref{fig:i1-h8-s}画的是相位变化的幅值，斜率比净绕一周的构形增大三倍。负电荷只改变斜率的方向，不改变周期大小。

在任一单值规范变换下，两条路径各自增加同样的端点相位$q[\chi(D)-\chi(S)]/\hbar$，相减仍为零。因此增加的三倍灵敏度来自路径绕数，而不是改变粒子的电荷。若实验中只看到较小磁通周期，必须同时知道干涉路径的拓扑，不能单靠周期把载流子的电荷解释成$3e$。若部分轨迹穿过磁场区，还会出现经典轨道和其他相位的改变，本题的无场路径假设不再适用。
\end{solution}

\begin{exercise}[带有限电刷半径的同极圆盘]\label{ex:i1-h9}
厚度$d$、电阻率$\rho$的导电圆盘以角速度$\omega>0$绕正$z$轴转动，轴向匀强磁场为$B\vct e_z$，$B>0$。内外电刷分别接在半径$0<a<R$与$R$处，经电阻$R_0$闭合。求电动势、盘内电阻、电流与维持匀速的转矩，忽略自感。 见图~\ref{fig:i1-h9-q}。
\end{exercise}

\begin{figure}[H]\centering
\begin{tikzpicture}[>=Stealth,line cap=round,line join=round]
\fill[gray!10](0,0)circle(1.5);\draw[curve](0,0)circle(1.5);\draw[curve](0,0)circle(.45);\draw[line width=1pt](.45,0)--(1.5,0);\node[above]at(.7,0){径向电流};\node at(-.7,.6){$\odot B$};\draw[->](1.75,0)arc(0:60:1.75)node[above]{$\omega$};\node[below]at(.45,0){$a$};\node[below]at(1.5,0){$R$};\node[right]at(2,0){两电刷经$R_0$接通};
\end{tikzpicture}
\caption{盘内电流径向流动，外接导线与电刷不随盘转动。}\label{fig:i1-h9-q}
\end{figure}
\begin{solution}
图~\ref{fig:i1-h9-q}的圆盘材料速度为$\vct u=\omega r\vct e_\theta$，磁场为$B\vct e_z$。因此$\vct u\times\vct B=B\omega r\vct e_r$，正电荷受力朝外。内刷到外刷的动生电动势为
\[\boxed{\mathcal E=\int_a^RB\omega r\dd r
=B\omega\left[\frac{r^2}2\right]_a^R
=\frac{B\omega}{2}(R^2-a^2)}.
\]
图~\ref{fig:i1-h9-s}中的梯形面积给出这个积分。电刷固定并不使材料速度为零，几何外电路的面积不变也不能抹去该项。


\begin{figure}[H]\centering
\begin{tikzpicture}[>=Stealth,line cap=round,line join=round]
\draw[axis](0,0)--(3.7,0)node[right]{$r$};\draw[axis](0,0)--(0,2.3)node[above]{$B\omega r$};\fill[gray!12](.6,0)--(.6,.36)--(3,1.8)--(3,0)--cycle;\draw[curve](0,0)--(3.3,1.98);\draw[guide](.6,0)--(.6,.36);\draw[guide](3,0)--(3,1.8);\node[below]at(.6,0){$a$};\node[below]at(3,0){$R$};\node[right]at(1.6,.4){面积为$\mathcal E$};
\end{tikzpicture}
\caption{洛伦兹力沿径向的积分，与盘面是否围成固定回路无关。}\label{fig:i1-h9-s}
\end{figure}

圆盘中半径$r$的环带截面为$2\pi rd$，径向电阻元为$\dd R=\rho\dd r/(2\pi rd)$，故
\[R_{\rm disk}=\frac\rho{2\pi d}\int_a^R\frac{\dd r}r
=\frac\rho{2\pi d}[\ln r]_a^R=\frac\rho{2\pi d}\ln\frac Ra,\qquad
I=\frac{\mathcal E}{R_0+R_{\rm disk}}.
\]
电流在盘内沿径向向外。磁力对径向电流元的转矩为$\dd\tau_z=-IBr\dd r$，于是所需驱动转矩大小
\[\begin{aligned}
\tau_{\rm drive}&=IB\int_a^Rr\dd r
=IB[r^2/2]_a^R=\frac{IB}{2}(R^2-a^2),\\
\tau_{\rm drive}\omega&=I\mathcal E=I^2(R_0+R_{\rm disk}).
\end{aligned}\]
驱动转矩与转动方向相同，外界输入的机械功全部转化为焦耳热。理想点接触$a\to0$时，盘内电阻对数发散；若要保持有限电流，不能同时把中心电刷面积取零又忽略接触附近的电阻。
\end{solution}

\begin{exercise}[滚轮上三个端点的电动势]\label{ex:i1-h10}
沿用双半径滚轮，外半径$b$、内半径$a<b$，顺时针纯滚动，磁场朝纸内。除最低点$A$与内圆最高点$B$外，再取外圆最高点$D$。求$A\to D$与$B\to D$的动生电动势。 见图~\ref{fig:i1-h10-q}。
\end{exercise}

\begin{figure}[H]\centering
\begin{tikzpicture}[>=Stealth,line cap=round,line join=round]
\draw[curve](0,0)circle(1.45);\draw[curve](0,0)circle(.75);\draw[curve](-2,-1.45)--(2,-1.45);\fill(0,-1.45)circle(1.4pt)node[below]{$A$};\fill(0,.75)circle(1.4pt)node[left]{$B$};\fill(0,1.45)circle(1.4pt)node[above]{$D$};\draw[->](0,0)--(.9,0)node[below]{$\omega b$};\node at(-.9,-.45){$\otimes\,\vct B$};
\end{tikzpicture}
\caption{三个端点距瞬心的距离不同。}\label{fig:i1-h10-q}
\end{figure}
\begin{solution}
图~\ref{fig:i1-h10-q}中，$A$仍为瞬心；$B,D$到瞬心的距离分别为$a+b$与$2b$。取磁场朝纸内、轮顺时针纯滚动，动生场可写成$\nabla\chi$，其中$\chi=B\omega[x^2+(y+b)^2]/2$。


\begin{figure}[H]\centering
\begin{tikzpicture}[>=Stealth,line cap=round,line join=round]
\draw[axis](0,-1.3)--(0,1.8)node[above]{$y$};\draw[axis](0,-1.3)--(3,-1.3)node[right]{$\chi$};\draw[curve,domain=-1.3:1.4,samples=81]plot({.32*(\x+1.3)^2},\x);\node[left]at(0,-1.3){$A$};\draw[guide](0,.7)--(1.28,.7);\draw[guide](0,1.3)--(2.163,1.3);\node[left]at(0,.7){$B$};\node[left]at(0,1.3){$D$};\node[right]at(2.2,1.4){$\chi=B\omega(y+b)^2/2$};
\end{tikzpicture}
\caption{辅助函数的端点差可避免重复计算曲线路径。}\label{fig:i1-h10-s}
\end{figure}

图~\ref{fig:i1-h10-s}在竖直直径上画出该辅助函数。沿任何导电路径的开放积分等于端点函数值之差，所以
\[\mathcal E_{A\to D}=2B\omega b^2,
\qquad
\boxed{\mathcal E_{B\to D}=\frac{B\omega}{2}\bigl[4b^2-(a+b)^2\bigr]}.
\]
由于$a<b$，后一个值为正，即正电荷沿轮内从$B$向$D$受到正的动生驱动。也可先算$A\to D$再减$A\to B$；开放线积分的方向必须随路径相减一起改变。

这里的$\chi$是动生项$\vct u\times\vct B$的积分辅助函数；整个系统的总电场还取决于电荷分布与电路状态。若电路开路并达到导电平衡，表面电荷会建立相应的电场；若有电流，盘内还存在欧姆压降。改变电刷接点并不等于改变轮的转速，计算时应重新代入端点到瞬心的距离，不能只把原答案中的一段几何长度随意替换。
\end{solution}

\begin{exercise}[无阻滚柱的非零初始电流]\label{ex:i1-h11}
均匀实心圆柱内固定面积$S$、自感$L$的无阻线圈，在水平面上纯滚动。初始法向平行竖直磁场$B$，但回路已有电流$I_0$。求转角$\psi=0$附近的稳定性，并在可作线性小振动时求频率。 见图~\ref{fig:i1-h11-q}。
\end{exercise}

\begin{figure}[H]\centering
\begin{tikzpicture}[>=Stealth,line cap=round,line join=round]
\draw[curve](-2,-1)--(2,-1);\draw[curve](0,0)circle(1);\draw[->](0,0)--(0,1.8)node[above]{$\vct B$};\draw[guide](-.8,-.4)--(.8,.4);\node[right]at(1.2,.5){$I(0)=I_0$};\draw[->](1.25,.8)arc(35:-30:1.4)node[right]{$\psi$};
\end{tikzpicture}
\caption{正I0对应初始磁矩与外磁场同向。}\label{fig:i1-h11-q}
\end{figure}
\begin{solution}
图~\ref{fig:i1-h11-q}的圆柱放在水平面上，因此没有沿坡重力项。初始磁链是$LI_0+BS$，不是$BS$。转角为$\psi$时，
\[LI+BS\cos\psi=LI_0+BS,\qquad
I=I_0+\frac{BS}{L}(1-\cos\psi).
\]
减去初始磁能后，有效势为
\[U-U_0=\frac L2\left[I_0+\frac{BS}{L}(1-\cos\psi)\right]^2-\frac L2I_0^2.
\]

\begin{figure}[H]\centering
\begin{tikzpicture}[>=Stealth,line cap=round,line join=round]
\draw[axis](-2,0)--(2.1,0)node[right]{$\psi$};\draw[axis](0,-1.1)--(0,2.1)node[above]{$U-U_0$};\draw[curve,domain=-1.3:1.3]plot(\x,{.7*(\x)^2});\draw[guide,domain=-1.3:1.3]plot(\x,{-.5*(\x)^2});\draw[densely dotted,domain=-1.3:1.3]plot(\x,{.3*(\x)^4});\node[right]at(1.3,1.3){$I_0>0$};\node[right]at(1.3,-.85){$I_0<0$};\node[right]at(1.4,.5){$I_0=0$};
\end{tikzpicture}
\caption{三种初始电流分别给出正二次、负二次或正四次首项。}\label{fig:i1-h11-s}
\end{figure}

图~\ref{fig:i1-h11-s}中的三种局部形状，来自$1-\cos\psi=\psi^2/2+O(\psi^4)$，故
\[U-U_0=\frac{I_0BS}{2}\psi^2+O(\psi^4).
\]
有效惯量由平动与转动动能相加得到：
\[
T=\frac12m(r\dot\psi)^2+\frac12\left(\frac12mr^2\right)\dot\psi^2
=\frac12J_{\rm eff}\dot\psi^2,\qquad J_{\rm eff}=\frac32mr^2.
\]
记$I(\psi)=I_0+(BS/L)(1-\cos\psi)$。链式法则给$I'(\psi)=(BS/L)\sin\psi$，故
\[\frac{\dd U}{\dd\psi}=LI(\psi)I'(\psi)
=BS\left[I_0+\frac{BS}{L}(1-\cos\psi)\right]\sin\psi.\]
代入$J_{\rm eff}\ddot\psi=-U'(\psi)$，原运动方程为
\[
J_{\rm eff}\ddot\psi
=-BS\left[I_0+\frac{BS}{L}(1-\cos\psi)\right]\sin\psi.
\]
保留一阶项后是$J_{\rm eff}\ddot\psi=-I_0BS\psi$。

对$I_0>0$，二阶系数为正，$J_{\rm eff}=3mr^2/2$给出小振动角频率
\[\boxed{\omega_0^2=\frac{I_0BS}{J_{\rm eff}}=\frac{2I_0BS}{3mr^2}}.
\]
对$I_0<0$，该构形不稳定。对$I_0=0$，二阶项消失，但首项为$B^2S^2\psi^4/(8L)>0$，因此仍是非线性稳定的平衡；振动频率依赖振幅，不能把二阶频率为零解释成完全没有回复作用。式中$I_0$按初始法向与磁场同向为正，若改变电流符号约定，稳定性判断须随物理磁矩一起转换。
\end{solution}

\begin{exercise}[圆磁区外沿螺线运动的功]\label{ex:i1-h12}
圆磁区半径$R$，区内$\partial_tB_z=-\alpha$、区外$B=0$。正电荷$q$、质量$m$的珠子从静止沿光滑绝缘螺线$r(\theta)=2R\exp[(\ln2)\theta/\pi]$运动，$0\le\theta\le\pi$。求电场作功与末速，并说明多绕一周的变化。 见图~\ref{fig:i1-h12-q}。
\end{exercise}

\begin{figure}[H]\centering
\begin{tikzpicture}[>=Stealth,line cap=round,line join=round]
\draw[guide](0,0)circle(.55);\fill[gray!10](0,0)circle(.55);\draw[curve,domain=0:180,samples=101]plot({1.1*exp(ln(2)*\x/180)*cos(\x)},{1.1*exp(ln(2)*\x/180)*sin(\x)});\fill(1.1,0)circle(1.5pt)node[right]{$A$};\fill(-2.2,0)circle(1.5pt)node[left]{$D$};\node at(0,0){$R$};\draw[axis](-2.5,0)--(1.7,0);
\end{tikzpicture}
\caption{轨迹半径从2R增至4R，始终在磁区外。}\label{fig:i1-h12-q}
\end{figure}
\begin{solution}
图~\ref{fig:i1-h12-q}的螺线全在圆磁区外，故$E_\theta=\alpha R^2/(2r)$、$E_r=0$。沿任意平面曲线，线元可分成$\dd\vct l=\vct e_r\dd r+\vct e_\theta r\dd\theta$，因此
\[\dd W=q\vct E\cdot\dd\vct l
=\frac{q\alpha R^2}{2}\dd\theta.
\]
半径$r(\theta)$完全消去。沿图中路径由$0$积到$\pi$，
\[W=\frac{q\alpha R^2}{2}\int_0^\pi\dd\theta
=\frac{q\alpha R^2}{2}[\theta]_0^\pi.\]
再用$mv_D^2/2=W$，得到
\[\boxed{W=\frac{q\alpha R^2\pi}{2}},\qquad
\boxed{v_D=\sqrt{\frac{q\alpha R^2\pi}{m}}}.
\]
第二式使用初始静止、光滑约束不做功，且外区没有其他做功力。图~\ref{fig:i1-h12-s}中功随展开极角线性增长，和螺线在外侧绕得更远并不矛盾：电场随$r$减弱，角向线元却随$r$增大。


\begin{figure}[H]\centering
\begin{tikzpicture}[>=Stealth,line cap=round,line join=round]
\draw[axis](0,0)--(4.2,0)node[right]{$\theta$};\draw[axis](0,0)--(0,2.1)node[above]{$W$};\draw[curve](0,0)--(3.5,1.65);\draw[guide](3.5,0)--(3.5,1.65);\node[below]at(3.5,0){$\pi$};\node[right]at(3.5,1.65){$q\alpha R^2\pi/2$};
\end{tikzpicture}
\caption{在给定场中，功由展开的极角决定。}\label{fig:i1-h12-s}
\end{figure}

若将同一螺线再延长一周到$\theta=3\pi$，终点半径为$2R\exp(3\ln2)=16R$，不会回到原终点。新增功为
\[\Delta W=\frac{q\alpha R^2}{2}\int_\pi^{3\pi}\dd\theta
=\frac{q\alpha R^2}{2}[\theta]_\pi^{3\pi}
=q\alpha\pi R^2.\]
若另选一条始终位于磁区外的轨道，多绕一周后仍到达原终点，则展开极角也比原路径多$2\pi$，新增功相同。两种比较的终点几何不同，功的变化都由连续展开的极角确定。式中角度须用弧度。
\end{solution}

\begin{exercise}[在静态梯度磁场中平移的圆环]\label{ex:i1-h13}
半径$a$、电阻$R$的圆环以速度$V\vct e_x$在$xy$平面内平移，外场$B_z=B_0+Gx$不随时间变。忽略自感，求电动势、电流和沿运动方向的净磁力，并核对机械功率。 见图~\ref{fig:i1-h13-q}。
\end{exercise}

\begin{figure}[H]\centering
\begin{tikzpicture}[>=Stealth,line cap=round,line join=round]
\draw[curve](-.7,0)circle(1.05);\draw[guide](.1,0)circle(1.05);\draw[->](-.7,0)--(.4,0)node[above]{$V$};\draw[axis](-2.1,-1.5)--(2.4,-1.5)node[right]{$x$};\node[below]at(-.7,-1.5){$X$};\node at(1.7,.8){$B_z=B_0+Gx$};\draw[->](-.7,0)--++(15:1.05)node[midway,below]{$a$};
\end{tikzpicture}
\caption{磁通变化来自圆环位置改变，而非外场的时间偏导。}\label{fig:i1-h13-q}
\end{figure}
\begin{solution}
图~\ref{fig:i1-h13-q}的外场不随时间变化，但圆环中心$X$在运动。圆盘关于自身中心对称，$x-X$的面积积分为零，所以
\[\Phi=\int_SB_z\dd S=\pi a^2(B_0+GX),\qquad
\mathcal E=-\pi a^2GV.
\]
记$A=\pi a^2$，忽略自感时$I=-AGV/R$。该式按逆时针为正；$G,V>0$时，图~\ref{fig:i1-h13-s}中的电流为顺时针。


\begin{figure}[H]\centering
\begin{tikzpicture}[>=Stealth,line cap=round,line join=round]
\draw[curve](0,0)circle(1.1);\draw[->](1.1,0)arc(0:-70:1.1);\draw[->](0,-.3)--(-1,-.3)node[left]{$F_x$};\draw[->](0,.3)--(1,.3)node[right]{$V$};\node at(0,0){$I<0$};\node[below]at(0,-1.6){$-F_xV=I^2R$};
\end{tikzpicture}
\caption{梯度场中的感应电流产生与速度相反的净磁力。}\label{fig:i1-h13-s}
\end{figure}

净磁力可直接对导线积分。其$x$分量为
\[F_x=I\oint B_z\dd y
=IG\oint x\dd y=IGA,
\]
在圆周上取$x=X+a\cos\theta$、$y=Y+a\sin\theta$，则$\dd y=a\cos\theta\dd\theta$，所以
\[\begin{aligned}
\oint x\dd y
&=Xa\int_0^{2\pi}\cos\theta\dd\theta
+a^2\int_0^{2\pi}\cos^2\theta\dd\theta\\
&=Xa[\sin\theta]_0^{2\pi}
+a^2\left[\frac\theta2+\frac{\sin2\theta}4\right]_0^{2\pi}
=\pi a^2=A.
\end{aligned}\]
其中均匀场项$IB_0\oint\dd y=IB_0[y]_{\text{一周}}=0$。代入$I$，
\[\boxed{F_x=-\frac{A^2G^2}{R}V},\qquad
-F_xV=\frac{A^2G^2V^2}{R}=I^2R.
\]
线性磁场使面积平均恰等于中心值，故这里不需假设$a$极小；但整个环必须处于给定的线性场范围内。若$G=0$，匀强场中的刚性整体平移没有闭路电动势，净磁阻力也为零。局部导线上的动生力仍可能存在，它们的完整回路积分互相抵消。
\end{solution}

\begin{exercise}[旋度、散度与中心值是否足够]\label{ex:i1-h14}
在整个平面给定$\partial_tB_z=-\alpha$，且知道$E_z=0$、$\nabla\cdot\vct E=0$、$\vct E(O)=0$。这些条件能否唯一推出$\vct E=(\alpha/2)\vct e_z\times\vct r$？构造一族满足同样条件的场，并说明还需要何种物理信息。 见图~\ref{fig:i1-h14-q}。
\end{exercise}

\begin{figure}[H]\centering
\begin{tikzpicture}[>=Stealth,line cap=round,line join=round]
\draw[axis](-1.7,0)--(1.9,0)node[right]{$x$};\draw[axis](0,-1.5)--(0,1.7)node[above]{$y$};\node at(-.8,.8){$-\dot B_z=\alpha$};\fill(0,0)circle(1.5pt)node[below left]{$O$};\node[below]at(0,-1.9){已知：散度零、旋度固定、中心场零};
\end{tikzpicture}
\caption{局部方程与单点值并不等于完整边值问题。}\label{fig:i1-h14-q}
\end{figure}
\begin{solution}
图~\ref{fig:i1-h14-q}所列条件仍只给出局部方程与一个点值。先取一个解
\[\vct E_0=(\alpha/2)(-y,x,0).\]
任取常数$\eta$，函数$f=\eta xy$满足$\nabla^2f=0$，且原点梯度为零，因此
\[\vct E_\eta=\vct E_0-\nabla f
=\left[-\left(\frac\alpha2+\eta\right)y,
\left(\frac\alpha2-\eta\right)x,0\right]
\]
对所有$\eta$都满足$\nabla\cdot\vct E_\eta=0$、$\nabla\times\vct E_\eta=\alpha\vct e_z$和$\vct E_\eta(O)=0$。


\begin{figure}[H]\centering
\begin{tikzpicture}[>=Stealth,line cap=round,line join=round]
\begin{scope}[xshift=-2.4cm]\draw[axis](-1.2,0)--(1.3,0);\draw[line width=.4pt](0,-1.2)--(0,0);\draw[->](.9,0)--(.9,.5);\draw[->](-.9,0)--(-.9,-.5);\draw[->](0,.9)--(-.5,.9);\node[below]at(0,-1.5){纯旋流};\end{scope}
\begin{scope}[xshift=2.4cm]\draw[axis](-1.2,0)--(1.3,0);\draw[line width=.4pt](0,-1.2)--(0,0);\draw[->](.8,0)--(.8,-.5);\draw[->](0,.8)--(-.5,.8);\draw[->](-.8,0)--(-.8,.5);\node[below]at(0,-1.5){可叠加$-\nabla(\eta xy)$};\end{scope}
\end{tikzpicture}
\caption{调和梯度分量改变局部电场，但不改变给定旋度。}\label{fig:i1-h14-s}
\end{figure}

图~\ref{fig:i1-h14-s}把纯旋流与可添加的调和梯度场分开画出。它们的环量相同，但一般点的电场方向与大小不同，所以题给条件不足以唯一确定电场。

若再要求整个物理源及边界具有连续轴对称，$\eta$项会破坏这种对称，才被排除。在有限区域中，也可以用完整的电势或电场边界条件消除该自由度。只说“没有体电荷”并不等于没有外部电极或表面电荷产生的调和场。无限空间均匀变化磁场的场本来随距离增长，不能再把不相容的“远处电场为零”作为补充条件；应回到一个有限源模型说明远区行为。
\end{solution}

